\documentclass[11pt,letterpaper,reqno]{amsart}
% BEGIN INLINE manuscript-style
\makeatletter
\NeedsTeXFormat{LaTeX2e}[2020/10/01]
% Use with \documentclass[11pt,letterpaper,reqno]{amsart}.
% This file controls presentation only. It supplies no mathematical content.

\RequirePackage[T1]{fontenc}
\RequirePackage{lmodern}
\RequirePackage{microtype}
\RequirePackage[letterpaper,margin=1in,headheight=15pt]{geometry}
\RequirePackage{amsmath,amsthm,amssymb,mathtools}
\RequirePackage{enumitem}
\RequirePackage{perpage}
\RequirePackage{longtable}
\RequirePackage{xcolor}
\definecolor{manuscriptred}{rgb}{0.6,0.1,0.1}
\definecolor{manuscriptblue}{rgb}{0.1,0.3,0.5}
\definecolor{manuscriptmagenta}{HTML}{60445C}
\RequirePackage[breaklinks=true,colorlinks=true,
  linkcolor=manuscriptred,citecolor=manuscriptblue,urlcolor=manuscriptmagenta]{hyperref}
\RequirePackage{cleveref}
\hypersetup{pdfauthor={}}

% Retain the AMS paragraph, title, heading and running-head conventions.
% Equation counters restart each section: 2.i, 2.ii, ... .
% Appendix sections keep their letters: A.i, A.ii, ... .
\numberwithin{equation}{section}
\renewcommand{\theequation}{\thesection.\roman{equation}}
% Prose references include "eq." and omit parentheses around the number.
\renewcommand{\eqref}[1]{\textup{eq.~\ref{#1}}}
\crefname{equation}{eq.}{eqs.}
\Crefname{equation}{eq.}{eqs.}
\creflabelformat{equation}{#2#1#3}
\setcounter{tocdepth}{1}
\renewcommand{\thefigure}{\Roman{figure}}
\renewcommand{\thetable}{\Roman{table}}
\newcommand{\gineedspace}[1]{\par\begingroup
  \dimen0=\pagegoal\advance\dimen0 by-\pagetotal
  \ifdim\dimen0<#1\relax\newpage\fi\endgroup}
\MakePerPage{footnote}
\renewcommand{\thefootnote}{%
  \ifnum\value{footnote}>9\arabic{footnote}\else\fnsymbol{footnote}\fi}

% Upright statements. Main results are A, B, C; supporting statements
% share a section-local sequence. These end marks indicate typography,
% never proof completeness, validation, or a deferred-proof status.
\theoremstyle{definition}
\newtheorem{ms@theorem}{Theorem}
\renewcommand{\thems@theorem}{\Alph{ms@theorem}}
\newtheorem*{ms@theoremstar}{Theorem}
\newtheorem*{ms@theoremzero}{Theorem 0}
\newtheorem{ms@proposition}{Proposition}[section]
\newtheorem{ms@lemma}[ms@proposition]{Lemma}
\newtheorem{ms@definition}[ms@proposition]{Definition}
\newtheorem{ms@corollary}[ms@proposition]{Corollary}
\newtheorem*{ms@corollarystar}{Corollary}

\newcommand{\ms@majorstart}{%
  \pushQED{\qed}\renewcommand{\qedsymbol}{$\blacksquare\square$}}
\newcommand{\ms@minorstart}{%
  \pushQED{\qed}\renewcommand{\qedsymbol}{$\blacksquare$}}
\newenvironment{theorem}[1][]{\ms@majorstart\ms@theorem[#1]}
  {\popQED\endms@theorem}
\newenvironment{theorem*}[1][]{\ms@majorstart\ms@theoremstar[#1]}
  {\popQED\endms@theoremstar}
\newenvironment{theoremzero}[1][]{\ms@majorstart\ms@theoremzero[#1]}
  {\popQED\endms@theoremzero}
\newenvironment{proposition}[1][]{\ms@minorstart\ms@proposition[#1]}
  {\popQED\endms@proposition}
\newenvironment{lemma}[1][]{\ms@minorstart\ms@lemma[#1]}
  {\popQED\endms@lemma}
\newenvironment{definition}[1][]{\ms@minorstart\ms@definition[#1]}
  {\popQED\endms@definition}
\newenvironment{corollary}[1][]{\ms@majorstart\ms@corollary[#1]}
  {\popQED\endms@corollary}
\newenvironment{corollary*}[1][]{\ms@majorstart\ms@corollarystar[#1]}
  {\popQED\endms@corollarystar}
% The source's lemma/proposition variants are deliberately not carried over:
% choose one marker for each statement type rather than encode proof status.

\theoremstyle{remark}
\newtheorem{remark}[ms@proposition]{Remark}
\newtheorem*{remark*}{Remark}
% The ordinary amsthm proof environment retains its open square.

\crefname{ms@theorem}{theorem}{theorems}
\Crefname{ms@theorem}{Theorem}{Theorems}
\crefname{ms@proposition}{proposition}{propositions}
\Crefname{ms@proposition}{Proposition}{Propositions}
\crefname{ms@lemma}{lemma}{lemmas}
\Crefname{ms@lemma}{Lemma}{Lemmas}
\crefname{ms@definition}{definition}{definitions}
\Crefname{ms@definition}{Definition}{Definitions}
\crefname{ms@corollary}{corollary}{corollaries}
\Crefname{ms@corollary}{Corollary}{Corollaries}
\crefname{remark}{remark}{remarks}
\Crefname{remark}{Remark}{Remarks}

% Bibliography configuration for this manuscript.
\makeatother
% END INLINE manuscript-style


\title[Gaussian block integration]{Gaussian block integration and free action-density correlations}
\date{Working draft, 9 November 2026}
\begin{document}
\raggedbottom
\begin{abstract}
We derive Gaussian integration on prescribed linear averaging fibres,
with physical measures and residual gauge factors. The constrained
minimizer, fluctuation covariance and effective action compose under
successive block integrations. Fixed real coordinate integrals admit
complex covariance parameters with coherent determinant branches,
local moment bounds and a cost determined by the selected coordinates.
For the flat four-dimensional curvature we compute the separated
action-density correlations and compare their lattice and continuum
kernels on a uniform class of separated tests. The quadratic data are
given; nonlinear background admission and interacting comparisons are
separate inputs.
\end{abstract}
\maketitle
\tableofcontents

\section{Introduction}

Block integration separates a field into the minimizer with its
prescribed average and a fluctuation with zero average. At the
quadratic level this is an exact Gaussian decomposition. Its
normalization depends on the constraint measure and, in a gauge
theory, on the coordinates used along residual gauge directions.
Those factors must be retained when integrations are composed or a
gauge is changed.

We first prove the finite-dimensional constrained formula and its
composition law. The proof needs positivity on the constrained
space, not on every direction of the ambient field space. We then
identify the coarea and gauge determinants in physical coordinates.
The complex integral is defined on fixed real coordinates, with a
determinant logarithm continued from a positive real form. This
permits fixed measurable cutoffs and gives estimates depending on
the selected dimension. A real Gaussian restricted by convex
coordinate tests has separate conditional and flux identities.

The last part concerns the free curvature in four dimensions.
Its separated Wick contractions have an explicit positive
two-point coefficient. We keep the finite-lattice contact term,
the periodic zero mode and the errors from mesh, volume and
coarse conditioning. Continuum cancellations are stated at their
actual scope; they are not exact assertions about every finite
lattice. The continuum Wick-square notation describes separated
correlation limits, not a claimed ultraviolet limit of each
individual random observable.

The Gaussian block transform and its constrained shift occur in
\cite[\S\S A--B, D]{Bal84a}; mixed-scale constraints are treated
in \cite[\S A]{Bal84b}. The background-field covariance is
constructed in \cite[\S E, Theorem~3.15]{Bal85p}. We use the
averaging geometry of \cite{BA} and the specified linear inverse,
coercivity and local matrix estimates of \cite{MB}. The proofs
below give the finite measure identities and free-density
comparisons in the conventions stated here. They do not identify
an interacting measure with the free model.

In a physical application the complete constrained Hessian includes
the multiplier paired with the second derivative of the constraint.
Its identification and the admission of a nonlinear critical
background must be supplied. Nonquadratic action, nonconstant
Haar and chart factors, and interacting cutoff corrections are
outside the present Gaussian calculation.

\section{A Gaussian integral with a linear constraint}
\label{gi:sec-real-gaussian}

Let \(E,F\) be finite-dimensional real Hilbert spaces. We use
the Lebesgue measure of their orthonormal coordinates, including
the physical weights when the spaces are lattice fields. The
adjoint for these pairings is \(Q^*\). On a complexification
we continue it by bilinear transpose. Determinants on the zero
space and the mass of its Lebesgue measure are one.

For $\operatorname{SU}(2)$ we use traceless Hermitian Lie
coordinates with real pairing
$\langle X,Y\rangle=\tfrac12\operatorname{tr}(XY)$.
The Pauli matrices are then an orthonormal real basis.

Let \(Q:E\to F\) be onto and put \(K=\ker Q\), with its
induced inner product and inclusion \(\iota_K:K\to E\).
Write \(J_Q=(\det QQ^*)^{1/2}\). For every \(b\in F\)
define the measure \(d\nu_{Q,b}\) on \(Q^{-1}(b)\) as
\(J_Q^{-1}\) times its induced Lebesgue measure. We use
\(\delta_F(Qx-b)\,dx\) for this measure. Indeed the
orthogonal right inverse \(Q^*(QQ^*)^{-1}\) maps \(F\)
to \(K^\perp\) with volume factor \(J_Q^{-1}\). This
proves the coarea formula with \(db\) the declared Lebesgue
measure of \(F\). Replacing that right inverse by any
other right inverse adds a shear into \(K\) and preserves
the volume factor. If \(S:F\to F'\) is an isomorphism,
then \(d\nu_{SQ,Sb}=|\det S|^{-1}d\nu_{Q,b}\), since the
fibre is the same and \(J_{SQ}=|\det S|J_Q\).

For a symmetric precision \(T:E\to E\) put \(T_K=\iota_K^*T\iota_K\).
If \(T_K>0\), write
\[
 I(T,Q;b)=\int_E e^{-\langle x,Tx\rangle/2}\,d\nu_{Q,b}(x),
 \qquad I(T,Q)=I(T,Q;0).
\]
The constrained shift and the normalization below are the finite
linear algebra of the block transform in \cite[\S D, eqs.~(1.64)--(1.65)]{Bal84a}.

\begin{theorem}[Constrained minimizer and fluctuation law]
\label{gi:thm-constrained-gaussian}
Suppose that \(Q:E\to F\) is onto, \(T=T^*\), and \(T_K>0\).
There is a unique linear map \(H:F\to E\) with
\[
 QH=I_F,\qquad \iota_K^*TH=0 .
\]
For \(Qx=b\), the minimum of \(\langle x,Tx\rangle/2\)
is attained uniquely at \(Hb\) and equals
\(\langle b,\Sigma b\rangle/2\), where \(\Sigma=H^*TH\).
Let \(C=\iota_KT_K^{-1}\iota_K^*\) and let \(\gamma_C\)
be the probability law on \(E\) of the centred Gaussian
on \(K\) with covariance \(T_K^{-1}\). Then for each
real \(b\) and bounded Borel \(\varphi:E\to\mathbb C\),
\begin{equation}\label{gi:eq-constrained-gaussian}
 \int_E \varphi(x)e^{-\langle x,Tx\rangle/2}\,d\nu_{Q,b}(x)
     = I(T,Q)e^{-\langle b,\Sigma b\rangle/2}
                            \int_E\varphi(Hb+u)\,d\gamma_C(u).
\end{equation}
The same identity holds for absolutely integrable \(\varphi\).
With \(m=\dim F\), the zero-fibre constant is
\begin{equation}\label{gi:eq-saddle-det}
 \begin{aligned}
 I(T,Q)&=(2\pi)^{\dim K/2}(\det T_K)^{-1/2}J_Q^{-1}
         =(2\pi)^{\dim K/2}|\det\mathbb S(T,Q)|^{-1/2},\\
 \det\mathbb S(T,Q)&=(-1)^m J_Q^2\det T_K,\qquad
 \mathbb S(T,Q)=
        \begin{pmatrix}T&Q^*\\Q&0\end{pmatrix}.
 \end{aligned}
\end{equation}
\end{theorem}

\begin{proof}
For a right inverse \(L_Q\) of \(Q\), the map
\[
 H=L_Q-\iota_K T_K^{-1}\iota_K^*TL_Q
\]
is still a right inverse and satisfies \(\iota_K^*TH=0\).
A difference of two such maps takes values in \(K\), where
the vanishing equation forces that difference to be zero.
Writing \(x=Hb+u\), \(u\in K\), removes the cross term in
the precision. The remaining quadratic form in \(u\) is
positive. This proves the minimizer and the effective
quadratic action. Coarea gives \(J_Q^{-1}du\) on the
fibre in these coordinates. Diagonalization of \(T_K\)
and the one-dimensional Gaussian integral give the first
line of \eqref{gi:eq-saddle-det} and the probability law
in \eqref{gi:eq-constrained-gaussian}.

Use the linear isomorphism
\((u,b,\lambda)\mapsto(\iota_Ku+Hb,\lambda)\)
on the two spaces of the saddle. Its absolute determinant is
\(J_Q^{-1}\). The saddle form in these coordinates is the
direct sum of \(T_K\) and
\(\begin{pmatrix}\Sigma&I_F\\I_F&0\end{pmatrix}\).
Exchanging the latter two block columns makes this matrix
block triangular with identity diagonal and contributes
the sign \((-1)^m\). This proves the determinant identity.
\end{proof}

The saddle is invertible and its inverse in the original
physical pairings is
\begin{equation}\label{gi:eq-saddle-inverse}
             \mathbb S(T,Q)^{-1}
                      =\begin{pmatrix}C&H\\H^*&-\Sigma\end{pmatrix}.
\end{equation}
In fact \(Tx+Q^*\lambda=f\), \(Qx=b\) first gives \(x=Cf+Hb\)
by testing on \(K\). Applying \(H^*\) to the first equation
then gives \(\lambda=H^*f-\Sigma b\). Adding a row quadratic
term \(Q^*A Q\), \(A=A^*\), leaves \(H,C,I(T,Q)\) unchanged
and replaces \(\Sigma\) by \(\Sigma+A\).

When \(T>0\) on all of \(E\), set \(G_0=T^{-1}\).
The equality \(TH=Q^*\Sigma\) follows from stationarity and
the fact that \(\operatorname{ran}Q^*=K^\perp\). Applying
\(QG_0\) gives
\begin{equation}\label{gi:eq-positive-schur}
 \Sigma=(QG_0Q^*)^{-1},\qquad
 H=G_0Q^*\Sigma,\qquad C=G_0-HQG_0 .
\end{equation}
For the last formula the right side takes values in \(K\)
and its product with \(\iota_K^*T\) is \(\iota_K^*\).
These two properties characterize \(C\). A Gaussian of covariance \(G_0\)
therefore decomposes as \(Hb+u\), where \(b\) has covariance
\(\Sigma^{-1}\), \(u\) has law \(\gamma_C\), and the two
variables are independent. This follows by integrating
\eqref{gi:eq-constrained-gaussian} in \(b\). The constrained
formula requires only the positive restriction \(T_K\).

\section{Successive constraints}
\label{gi:sec-successive}

Let \(M:E\to F\) and \(N:F\to G\) be onto. Here \(G\) is
a third real Hilbert space with its declared Lebesgue measure.
For \(T=T^*\) positive on \(\ker(NM)\), use \(H_M,C_M,\Sigma_M\)
for the operators of \cref{gi:thm-constrained-gaussian}
with \((T,M)\). For the second step use \(H_N,C_N,\Sigma_N\)
with \((\Sigma_M,N)\). These exist: \(T\) is positive on
\(\ker M\), and \(H_M b\) is a nonzero vector in \(\ker(NM)\)
for every \(b\ne0\) in \(\ker N\). Thus \(\Sigma_M\) is
positive on \(\ker N\).

\begin{proposition}[Gaussian composition]
\label{gi:prop-successive}
Under these hypotheses the operators for \((T,NM)\) satisfy
\begin{equation}\label{gi:eq-composed-gaussian-operators}
 H_{NM}=H_M H_N,\qquad
 C_{NM}=C_M+H_M C_N H_M^*,\qquad
 \Sigma_{NM}=\Sigma_N .
\end{equation}
For every \(c\in G\) and every bounded Borel test
\(\varphi:E\to\mathbb C\),
\begin{align}
 &\int_E\varphi(x)e^{-\langle x,Tx\rangle/2}\,d\nu_{NM,c}(x)\notag\\
 &\qquad=I(T,M)I(\Sigma_M,N;c)\,
        \mathbb E\bigl[\varphi(H_MH_Nc+U+H_M V)\bigr],
 \label{gi:eq-successive-tested}
\end{align}
where \(U,V\) are independent with laws \(\gamma_{C_M}\)
and \(\gamma_{C_N}\). In particular
\begin{equation}\label{gi:eq-successive-normalizers}
                     I(T,NM)=I(T,M)I(\Sigma_M,N).
\end{equation}
\end{proposition}

\begin{proof}
Write \(x=u+H_M(v+H_Nc)\), where \(u\in\ker M\)
and \(v\in\ker N\). With \(c\) free this is an
isomorphism onto \(E\) of absolute determinant
\(J_M^{-1}J_N^{-1}\), by the two coarea changes.
At \(c=0\) it maps the first two slots onto \(\ker(NM)\).
The last slot is a right inverse to \(NM\), so the
product \(J_M^{-1}J_N^{-1}\,du\,dv\) is exactly
\(d\nu_{NM,c}\) in these coordinates for every \(c\).
The precision in them is
\[
 \langle x,Tx\rangle
  =\langle u,Tu\rangle+\langle v,\Sigma_Mv\rangle
                                  +\langle c,\Sigma_Nc\rangle
\]
by the two stationarity equations. Thus \(u=v=0\) is
the unique minimum on each fibre and the two centred
Gaussian factors are independent. Applying the
single-fibre normalization to them gives
\eqref{gi:eq-successive-tested}. Computing the covariance
of \(U+H_MV\) proves the middle identity in
\eqref{gi:eq-composed-gaussian-operators}. Taking the
constant test and \(c=0\) proves the product.
\end{proof}

For any finite chain of onto constraints, assume positivity on
the kernel of their total composition. Use the composite of
the remaining maps as \(N\) in \cref{gi:prop-successive}.
The effective precision of the first step is then positive on
the kernel of that remaining composite. Induction gives the
successive law for all the stages. A second-stage
quadratic weight on \(Mx\) can be included by adding it
to \(\Sigma_M\). For example, suppose that the supplied
complete first-stage Hessian is \(B_1\) and its lift \(H\)
satisfies \(MH=I\) and \(\iota_{\ker M}^*B_1H=0\).
Suppose that the first-stage gauge coordinate \(G_1\)
satisfies \(G_1H=0\). If the supplied total form is
\[
 B_1+G_1^*G_1+M^*(A_2+G_2^*G_2)M,
\]
the first-step restricted precision is
\((B_1+G_1^*G_1)|_{\ker M}\) as a quadratic form.
The second-step precision is \(H^*B_1H+A_2+G_2^*G_2\).
Positivity on the respective integration kernels permits
\eqref{gi:eq-successive-normalizers}. This is the algebraic
composition for a prescribed gauge and Hessian splitting.

The corresponding identity for logarithms uses the same
reference at both stages. If the first factor is evaluated
at preceding data \(p(z)\) and the second at \(z\), define
each normalized logarithm by subtracting its value at
\(p(z_*)\) or \(z_*\), respectively. The total normalized
logarithm is then the sum of the first evaluated at
\(p(z)\) and the second at \(z\). This follows by applying
\eqref{gi:eq-successive-normalizers} at \(z\) and \(z_*\).
Any constant measure factor cancels in that difference.

\subsection{The flat curvature slices}
\label{gi:sec-free-slices}

Fix an odd blocking factor \(L\ge3\) and a nonnegative integer \(K\)
of levels. Use a periodic hierarchy in which each fine
period is divisible by \(L^K\) and has at least two sites
at level \(K\), as in the admitted geometry of \cite{BA,MB}.
Let \(\mathcal E=\operatorname{ran}\mathrm d\) be the exact
fine two-cochains on a compatible homogeneous periodic
lattice, with the counting Hilbert pairing on raw
plaquette fluxes. Write \(Q_j^{\rm pl}\) and \(Q_j^{\rm st}\)
for the composed plain and straight link rows at the identity
in raw link coordinates. Write \(\mathrm d_j\) for the
coboundary at level \(j\), with \(\mathrm d=\mathrm d_0\).
Let \(\mathcal R_j\) be the actual
level-\(j\) sheet average on \(\mathcal E\), with
\(\mathcal R_0=I\). In raw flux units, for a coarse plaquette
\(P\) based at \(y\), this average is
\[
 (\mathcal R_j F)(P)=L^{-4j}
   \sum_{z\in B^j(y)}\ \sum_{p\in\operatorname{Sq}^j(z;P)} F(p).
\]
Here \(B^j(y)\) is the level-\(j\) block of \(L^{4j}\)
fine sites, and \(\operatorname{Sq}^j(z;P)\) is the square
of \(L^{2j}\) oriented fine plaquettes parallel to \(P\),
translated to start at \(z\). Multiplicities are retained
under periodic projection. Put \(W_j=\mathcal E\ominus\ker\mathcal R_j\)
and denote by \(P_j\) the orthogonal projection onto \(W_j\),
extended by zero on \(\mathcal E^\perp\). In particular
\(P_0=P_{\rm ex}\), the exact-form projection.
The sheet composition and deterministic projection identities
are given in \cite[The deterministic curvature slices]{MB}.
For the hard constraints fix the internal coordinate spanning
tree rooted at the center of each block, as in
\cite[Kernel of a flat curvature slice]{MB}.

\gineedspace{18\baselineskip}
\begin{proposition}[Independent curvature slices]
\label{gi:prop-free-slices}
Let \(F_{\rm free}\) be the centred Gaussian with covariance
\(g^2P_{\rm ex}\), \(g>0\). Its conditional mean given
\(\mathcal R_jF_{\rm free}\) is \(P_jF_{\rm free}\), the
least-norm exact cochain with that datum. The decomposition
\begin{equation}\label{gi:eq-free-slices}
 F_{\rm free}=P_KF_{\rm free}+
                    \sum_{j=1}^K(P_{j-1}-P_j)F_{\rm free}
\end{equation}
has independent centred terms, with covariances \(g^2P_K\)
and \(g^2(P_{j-1}-P_j)\). At step \(j\), take the
intermediate link fields whose next straight block average
\(S_La\) is zero and whose prescribed internal spanning-tree links
are zero. Parametrize this subspace by its remaining
nonpivot variables. If \(B_j\) sends these variables to
their least-norm fine curvature, and \(A_j=B_j^*B_j\) is the restricted
unit-coupling step precision in those variables, the precision
at coupling \(g\) is \(g^{-2}A_j\). Its fluctuation covariance
pushes forward to
\[
            B_j(g^2 A_j^{-1})B_j^*
                        =g^2(P_{j-1}-P_j).
\]
\end{proposition}

\begin{proof}
Use the orthogonal split
\(\mathcal E=W_j\oplus\ker\mathcal R_j\).
The two Gaussian components are independent since the
law on \(\mathcal E\) has scalar covariance. The restriction
of \(\mathcal R_j\) to \(W_j\) is one-to-one and onto the
range of \(\mathcal R_j\), so the datum fixes the first
component. The second is centred. Orthogonality gives
the asserted minimum. At the identity, the plain and
straight link derivatives differ by a coarse gradient;
their coarse coboundaries therefore agree. The exact
flux formula in \cite[The flat derivative and its curl]{BA}
gives, for both raw link rows,
\begin{equation}\label{gi:eq-plain-straight-flux}
 \mathrm d_j Q_j^{\rm pl}=\mathrm d_j Q_j^{\rm st}
                         =\mathcal R_j\mathrm d .
\end{equation}
Thus observing either averaged link field modulo
closed cochains gives the same conditioning, since the
coarse coboundary identifies that quotient with the
exact coarse cochains.

The ranges \(W_j\) are nested by sheet composition.
The successive differences of their projections are
orthogonal projections with pairwise orthogonal ranges.
Their telescoping sum proves \eqref{gi:eq-free-slices},
and orthogonal coordinates of the scalar Gaussian
prove the stated independence and covariances. Finally
the hard-tree parametrization and flat slice result of
\cite[Kernel of a flat curvature slice]{MB}
identify \(B_j\) as an injection onto
\(W_{j-1}\cap W_j^\perp\). For an injective map \(B_j\),
the orthogonal projection onto its range is
\(B_j(B_j^*B_j)^{-1}B_j^*\), as follows by testing
the error against that range. This gives the lifted
fluctuation covariance in the same counting units.
\end{proof}

\begin{proposition}[The flat curvature law in two gauges]
\label{gi:prop-plain-straight-curvature}
On the compatible periodic hierarchy of \cref{gi:prop-free-slices}
at the identity connection, use the plain and straight rows
and the counting inner products defined there.
For a step $1\le j\le K$ put $A=\mathcal R_{j-1}|_{\mathcal E}$ and
$V=A A^*$, viewed on $\mathcal E_{j-1}=\operatorname{ran}\mathrm d_{j-1}$.
The stored real cochains carry counting inner products. Their
linear subspaces carry the induced inner products and Lebesgue
measures, which also determine all adjoints.
Use the specified flat quadratic energy
\[
  \frac{1}{2g^2}\big\langle \mathrm d_{j-1}a,
                     V^{-1}\mathrm d_{j-1}a\big\rangle
\]
on the intermediate links. Constrain the next plain row to zero
and add the soft Gaussian penalty
$\langle G_{\rm con}a,W_{\rm con}G_{\rm con}a\rangle/2$,
where $W_{\rm con}$ is any fixed positive definite matrix
on the contour coordinates. Here $G_{\rm con}a$ is the
arithmetic mean of the contour path sums from the block center to each noncenter
site. Alternatively, constrain the next straight row to zero and
set the internal spanning-tree links to zero in each block. The two
normalized laws have the same curvature pushforward. It has
precision form $g^{-2}V^{-1}$ restricted to the exact intermediate
curvatures annihilated by the next plaquette average. Its least-norm
lift to the bare lattice has covariance $g^2(P_{j-1}-P_j)$.
\end{proposition}

\begin{proof}
Use one colour. Both one-step link rows are onto. A right inverse assigns the
datum of each positive coarse bond to all fine bonds crossing its
literal outgoing block face on the periodic lift, and sets internal
bonds to zero. Every translated
straight segment crosses that face once and every internal contour
sum is zero. This construction works on the stored bond occurrences,
including across the periodic seam. Composing and taking curls
shows that $A$ is onto $\mathcal E_{j-1}$. The least-norm lift
of an intermediate curvature $h$ is $A^*V^{-1}h$. Its squared
norm is $\langle h,V^{-1}h\rangle$. This is the flat effective
curl energy used in \cite[Proposition~4.1 and Corollary~4.3]{MB}.

Write $q_{\rm pl}$ for evaluation on block centers and $q_{\rm st}$
for the arithmetic mean on each block. On a site gradient
$\mathrm d_{j-1}\lambda$ the next raw link rows give, respectively,
$\mathrm d_j q_{\rm pl}\lambda$ and $\mathrm d_j q_{\rm st}\lambda$.
The first identity follows by telescoping each full contour from
center to center. The second follows by telescoping the straight
segments and averaging their endpoints. Both site rows are onto.
Both link rows send a constant directional field $c$ to $Lc$.
On either periodic lattice a closed link field is a gradient plus a
constant directional field. For instance, at each nonzero Fourier
frequency the coboundary symbol is wedge multiplication by the
nonzero difference vector, whose kernel on links is its span.
The zero frequency is the constant directional part.

Let $h$ be exact with zero next plaquette average and choose any
$a$ with $\mathrm d_{j-1}a=h$. Its next link datum is closed by
\eqref{gi:eq-plain-straight-flux} for a single step. Subtract a
constant fine field to remove the constant part of that datum,
then subtract the gradient of a site lift to remove its gradient
part. This works for either row. Thus the curl maps both zero-datum
spaces onto the same curvature subspace. Its null space for either
row consists of gradients with $q_{\rm pl}\lambda=0$ or
$q_{\rm st}\lambda=0$, respectively. Indeed the constant part of a
closed field in the zero-datum space must vanish, and a constant
site value can be subtracted from $\lambda$ to impose the indicated
condition.

On the plain null space $G_{\rm con}\mathrm d_{j-1}\lambda$ lists the
noncenter values of $\lambda$, since the contour sum is
$\lambda(x)-\lambda(y)$ and $\lambda(y)=0$ at the center $y$.
It is a bijection. Choose a linear lift of $h$ to the plain
zero-datum space and add the null-space gradient. This linear
splitting has a constant Jacobian, and the row coarea factor
$J_Q^{-1}$ is constant as well. Changing the
residual variables to $G_{\rm con}a$ in the Gaussian integral has constant
Jacobian, and translation by the contour values of the lift leaves
the integral unchanged. The soft integral hence multiplies the
curvature density by a constant independent of $h$.

For the straight row, integrate the desired tree link values
along each internal tree to obtain a site function. Subtract its
block means and then subtract its gradient from $a$. The row stays
zero, and the resulting tree links vanish. A residual gradient
vanishing on the tree has a constant potential in each block, and
its zero block means force it to vanish. The hard-tree space thus
maps bijectively to the same curvature subspace with constant
Jacobian. In both cases normalization cancels the constant, giving
the stated restriction of the curvature precision. The least-norm
lift maps that subspace onto $W_{j-1}\cap W_j^\perp$ by sheet
composition. Its energy is the bare squared norm divided by
$2g^2$, which proves the last covariance identity.
\end{proof}

For the background current of the unit-coupling Wilson
quadratic form and $F\in\mathcal E$, write
$P_jF=(\mathcal R_j|_{\mathcal E})^*\phi$ for some coarse
plaquette field $\phi$. The restricted adjoint includes
$P_{\rm ex}$, and $\mathrm d^*P_{\rm ex}=\mathrm d^*$
because $\mathrm d$ takes values in $\mathcal E$. Taking the transpose of
\eqref{gi:eq-plain-straight-flux} gives
\begin{equation}\label{gi:eq-free-background-current}
     \mathrm d^* P_jF
          =(Q_j^{\rm st})^*\mathrm d_j^*\phi .
\end{equation}
Thus the current has the row-multiplier form even in the flat
quadratic theory. This is also the deterministic identity in
\cite[The deterministic curvature slices]{MB}.
Multiplying the quadratic action by $g^{-2}$ multiplies
this current by the same factor.

\section{Measure factors and changes of gauge}
\label{gi:sec-gauge-measures}

The Gaussian normalizer depends on the measures of both the field
and constraint spaces. A gauge change preserves the physical
quadratic form but changes the coordinates integrated along its
null directions. We retain these two effects separately.

\subsection{Physical coordinates}

All determinants in this section use orthonormal coordinates in
the declared real Hilbert spaces. Suppose that a stored coordinate
array on $E$ has inner product $x^{\mathsf t}W_Ey$, with
$W_E>0$, and that $F$ has Gram matrix $W_F>0$. Then
\begin{equation}\label{gi:eq-coordinate-measures}
 dx_E=(\det W_E)^{1/2}\,dx,\qquad
 \delta_F(y)=(\det W_F)^{-1/2}\delta_{\rm array}(y),\qquad
 Q^*=W_E^{-1}Q^{\mathsf t}W_F.
\end{equation}
Here $\mathsf t$ is the ordinary array transpose. The exponent
of the operator $T$ is $x^{\mathsf t}W_ETx/2$. For the already defined $I(T,Q)$,
\begin{equation}\label{gi:eq-array-integral}
 I(T,Q)=\left(\frac{\det W_E}{\det W_F}\right)^{1/2}
 \int_{\mathbb R^{\dim E}}\delta_{\rm array}(Qx)
                         e^{-x^{\mathsf t}W_ETx/2}\,dx.
\end{equation}
Indeed $W_E^{1/2}$ and $W_F^{1/2}$ carry the two physical
inner products to Euclidean ones. Their Jacobians give the
first two formulas, and the defining identity for an adjoint
gives the third. These changes include every lattice-volume
weight. Replacing $Q^*$ by $Q^{\mathsf t}$ before making this
change would alter both the covariance and the coarea factor.

\subsection{A relative gauge normalizer}

Let $E,F,\mathcal H,\mathcal K$ be finite-dimensional real
Hilbert spaces. Assume injective maps
$D:\mathcal H\to E$ and $d:\mathcal K\to F$, onto maps
$Q:E\to F$ and $R:\mathcal H\to\mathcal K$, and the identity
\begin{equation}\label{gi:eq-gauge-chain}
                         QD=dR.
\end{equation}
Let $B=B^*$, $BD=0$, and suppose that $B$ is nonnegative on
$\ker Q$, with null space there exactly $D(\ker R)$.
Put $\mathcal H_0=\ker R$ and $D_0=D|_{\mathcal H_0}$.
A residual gauge coordinate is a map $G:E\to Z$, where
$Z$ is a real Hilbert space with a fixed norm and measure,
$\dim Z=\dim\mathcal H_0$, and
$GD_0:\mathcal H_0\to Z$ is invertible. Its determinant
uses the induced Hilbert measure of $\mathcal H_0$ and
the declared measure of $Z$.
The operator $B+G^*G$ is positive on $\ker Q$: a zero vector
for its quadratic form must equal $D_0v$, and $GD_0v=0$
then gives $v=0$.

\begin{theorem}[Relative gauge normalization]
\label{gi:thm-relative-gauge}
Under the preceding hypotheses, the quantity
\begin{equation}\label{gi:eq-relative-gauge-normalizer}
 \mathcal J(B,Q,R,G)
   =\sqrt{\det(RR^*)}\,|\det(GD_0)|\,I(B+G^*G,Q)
\end{equation}
does not depend on the residual coordinate $G$.
Let $S:E\to\mathcal K$, and suppose
\[
                 Q'=Q+dS,\qquad R'=R+SD
\]
are onto. If the same hypotheses hold for $Q',R'$ and an
admitted residual coordinate $G'$, then
\[
             \mathcal J(B,Q',R',G')=\mathcal J(B,Q,R,G).
\]
Only the coarse gauge component of the constraint changes.
\end{theorem}

\begin{proof}
We express the normalizer on the fixed physical quotient.
Use the orthogonal complements
$E_\perp=(\operatorname{ran}D)^\perp$ and
$F_\perp=(\operatorname{ran}d)^\perp$, and let $P$ be
the orthogonal projection onto $F_\perp$. Define
$\overline Q=PQ|_{E_\perp}$ and
$B_\perp=B|_{E_\perp}$ as a quadratic form. The map
$\overline Q$ is onto: for a vector of $F_\perp$, choose
a preimage under $Q$ and discard its $\operatorname{ran}D$
component, which changes $Q$ only in $\operatorname{ran}d$.
The restriction of $B_\perp$ to $\ker\overline Q$ is
positive. In fact $Qu=d\eta$ there, and an $h$ with
$Rh=\eta$ gives $u-Dh\in\ker Q$. Its action is that
of $u$. If it is zero, $u-Dh=D_0v$, whence $u$ lies in
$\operatorname{ran}D\cap E_\perp=\{0\}$.

Write $a=u+Dv$, with $u\in E_\perp$, $v\in\mathcal H$.
Its Jacobian is $J_D=\sqrt{\det(D^*D)}$. Set
$L=(d^*d)^{-1}d^*Q|_{E_\perp}$. Splitting $F$ into
$F_\perp\oplus\operatorname{ran}d$ gives
\[
 Qa=\overline Qu+d(Lu+Rv),\qquad
 \delta_F(Qa)=J_d^{-1}\delta_{F_\perp}(\overline Qu)
                                  \delta_{\mathcal K}(Lu+Rv),
 \qquad J_d=\sqrt{\det(d^*d)}.
\]
The coarea formula for $R$ now writes
$v=v_0-R^*(RR^*)^{-1}Lu$, $v_0\in\mathcal H_0$, with
factor $J_R^{-1}=(\det RR^*)^{-1/2}$. Since $BD=0$,
the action is $\langle u,Bu\rangle/2$. The gauge integral is
\[
 \int_{\mathcal H_0}
   \exp\{-\tfrac12|G(u-DR^*(RR^*)^{-1}Lu)+GD_0v_0|^2\}\,dv_0
       =\frac{(2\pi)^{(\dim\mathcal H-\dim\mathcal K)/2}}
                    {|\det(GD_0)|}.
\]
Combining these exact changes of variables proves
\begin{equation}\label{gi:eq-quotient-normalizer}
 \mathcal J(B,Q,R,G)=
 \frac{J_D}{J_d}(2\pi)^{(\dim\mathcal H-\dim\mathcal K)/2}
       I(B_\perp,\overline Q).
\end{equation}
Every integral is finite by the established quotient positivity.
The right side is independent of $G$. It is also unchanged by
$Q\mapsto Q+dS$, because $P d=0$; all its other data are fixed.
This proves both conclusions.
\end{proof}

\begin{corollary}[A residual gauge with a stabilizer]
\label{gi:cor-relative-gauge-stabilizer}
In \cref{gi:thm-relative-gauge}, allow \(D,d\) to have kernels.
Require \(D|_{\ker R}\) to be injective and keep the other
hypotheses, including that the null space of \(B\) on \(\ker Q\)
is \(D(\ker R)\). Then the same invariant
\(\mathcal J(B,Q,R,G)\), using the original Hilbert measures in
\eqref{gi:eq-relative-gauge-normalizer}, is independent of \(G\)
and is unchanged by the stated change of constraint. Require
the corresponding hypotheses for the new \(R',Q',G'\).
\end{corollary}

\begin{proof}
Write \(\mathcal H_{\rm s}=\ker D\),
\(\mathcal K_{\rm s}=\ker d\), and let
\(\mathcal H_{\rm e},\mathcal K_{\rm e}\) be their orthogonal
complements. The restriction \(A=R|_{\mathcal H_{\rm s}}\)
takes values in \(\mathcal K_{\rm s}\) by the gauge-chain identity.
It is injective by injectivity of \(D|_{\ker R}\).
If \(\zeta\in\mathcal K_{\rm s}\), choose \(v\) with \(Rv=\zeta\).
Then \(Dv\in\ker Q\) and \(BDv=0\). By the prescribed null
space there is \(v_0\in\ker R\) with \(Dv=Dv_0\).
Thus \(v-v_0\in\mathcal H_{\rm s}\) and \(A(v-v_0)=\zeta\).
This proves that \(A\) is an isomorphism.

In the orthogonal decompositions of the two site spaces, write
\[
 R=\begin{pmatrix}A&B_R\\0&C_R\end{pmatrix}.
\]
The map \(C_R:\mathcal H_{\rm e}\to\mathcal K_{\rm e}\)
is onto. The effective gradients \(D_{\rm e}=D|_{\mathcal H_{\rm e}}\)
and \(d_{\rm e}=d|_{\mathcal K_{\rm e}}\) are injective and
satisfy \(QD_{\rm e}=d_{\rm e}C_R\). The original \(\ker R\)
is the graph \(w\mapsto(-A^{-1}B_Rw,w)\) over \(\ker C_R\).
Let \(\pi:\ker R\to\ker C_R\) be the projection onto its
second component, with absolute determinant \(J_\pi\) in
the induced Hilbert measures. In particular
\(D|_{\ker R}=(D_{\rm e}|_{\ker C_R})\pi\).
The effective chain therefore has the same residual gauge
image and the same required nondegeneracy on \(\ker Q\).

The two coarea factors satisfy
\begin{equation}\label{gi:eq-stabilizer-coarea}
                 J_R J_\pi=|\det A|J_{C_R}.
\end{equation}
To verify this, use \(w\in\ker C_R\) and
\(r=(r_{\rm s},r_{\rm e})\in\mathcal K\) as coordinates
on \(\mathcal H\). Set
\[
 v_{\rm e}=w+L_C r_{\rm e},\qquad
 v_{\rm s}=A^{-1}(r_{\rm s}-B_Rv_{\rm e}),
\]
where \(L_C\) is the orthogonal right inverse of \(C_R\).
Then \(R(v_{\rm s},v_{\rm e})=r\), and the triangular
Jacobian of this change is \(|\det A|^{-1}J_{C_R}^{-1}\).
Alternatively, first use the coarea factor \(J_R^{-1}\)
with the induced measure on \(\ker R\), and then use
\(J_\pi^{-1}\) to express that measure in \(w\).
Equality of these Jacobians proves \eqref{gi:eq-stabilizer-coarea}.

Because \(|\det(GD_0)|=|\det(GD_{\rm e}|_{\ker C_R})|J_\pi\),
the original invariant equals \(|\det A|\) times
the invariant for the effective chain, with its induced
orthogonal Hilbert measures. Theorem~\ref{gi:thm-relative-gauge}
applies to the effective chain. Under \(R'=R+SD\) the map
\(A\) is unchanged, and the effective map is
\(C_R'=C_R+P_{\mathcal K_{\rm e}}SD_{\rm e}\).
This is exactly the same constraint change for that chain
because \(d\) annihilates \(\mathcal K_{\rm s}\).
The effective invariance and the fixed factor \(|\det A|\)
prove both claims with the original measures.
\end{proof}

The factor involving $R$ cannot be dropped. For example, take
$E=\mathbb R^2$, $F=\mathcal H=\mathcal K=\mathbb R$,
$Dv=(0,v)$, $d=1$, $Q(u,v)=\alpha u+\beta v$ with
$\beta\ne0$, and $B=\operatorname{diag}(\lambda,0)$,
$\lambda>0$. There is no residual gauge coordinate.
Then $R=\beta$ and
\[
 I(B,Q)=|\beta|^{-1}\sqrt{2\pi/\lambda},\qquad
 \mathcal J(B,Q,R,0)=\sqrt{2\pi/\lambda}.
\]
Changing the coarse gauge constraint changes the uncorrected
integral even though it leaves the physical quotient unchanged.

\subsection{The background gauge and stabilizers}

Set $\Delta=D^*D$, $P_0=I-R^*(RR^*)^{-1}R$, and use
$G=P_0D^*:E\to\mathcal H_0$. Then
$GD_0=\Delta_{00}$, the compression of $\Delta$ to
$\mathcal H_0$, and $G^*G=DP_0D^*$. The saddle determinant
formula for $\mathbb S(A,Q)=\left(\begin{smallmatrix}A&Q^*\\Q&0\end{smallmatrix}\right)$
therefore gives
\begin{equation}\label{gi:eq-background-gauge-normalizer}
 \mathcal J=
 \frac{|\det\mathbb S(\Delta,R)|}{\sqrt{\det(RR^*)}}
                    I(B+DP_0D^*,Q).
\end{equation}
Indeed $|\det\mathbb S(\Delta,R)|=
\det(RR^*)\det\Delta_{00}$; substitution into
\eqref{gi:eq-relative-gauge-normalizer} proves the formula.
When $RR^*=I$, its logarithm is
\begin{align}
 \log\mathcal J={}&
  -\tfrac12\log|\det\mathbb S(B+D(I-R^*R)D^*,Q)|
  +\log|\det\mathbb S(\Delta,R)|
  +\tfrac{\dim E-\dim F}{2}\log(2\pi).
 \label{gi:eq-background-gauge-log}
\end{align}
The second determinant is a constrained site determinant.
It is not an arbitrary determinant of a gauge operator.

For complex data all transposes are bilinear continuations
of the physical real adjoints. For a saddle with target dimension $m$, the real positive factor
is $(-1)^m\det\mathbb S(T,Q)$. Continue its logarithm and
square root from that positive value on the admitted branch
domain. The site saddle uses its own target dimension
$\dim\mathcal K$. Absolute values are not holomorphic
normalizations.

\subsection{The specified physical fibre}

We now state the use of the preceding algebra on lattice data.
The geometric maps and their derivatives are those of
\cite{BA}; the nondegenerate quadratic forms are the inputs
treated in \cite{MB}. No existence of a nonlinear critical
background is inferred from the following Gaussian identity.

Consider a homogeneous blocking with fine mesh $\eta$ and
coarse mesh one. Give fine site coordinates weight $\eta^4$
and coarse site coordinates weight one. Let $R_0$ evaluate
a site field at the distinct coarse nodes. Its physical
adjoint inserts $\eta^{-4}$ times each node value, so
\begin{equation}\label{gi:eq-node-coarea}
                         R_0R_0^*=\eta^{-4}I.
\end{equation}
For an assigned block with $N$ fine sites, the relative
transported average has the form
$R\lambda=N^{-1}\sum_x O_x\lambda_x$, where every $O_x$
is a real adjoint transport. The fine-to-coarse measure
ratio is $N^{-1}$. Therefore its physical adjoint has
components $O_x^*$, and
\begin{equation}\label{gi:eq-site-coisometry}
                       RR^*=N^{-1}\sum_x O_xO_x^*=I.
\end{equation}
Different blocks are disjoint. A homogeneous composition
of these coisometries is again a coisometry, by applying
$R_1R_1^*=I$ successively from the inside. The maps here
are relative derivatives centered at the given real
background. A derivative of a fixed-reference nonlinear
chart at a nonzero coordinate is a different map.

Let $T_{yx}$ be the mean contour transport from a parent
node $y$ to a site $x$ at the real reference background.
Endpoint covariance gives the derivative of its relative
left logarithm as
$\lambda(y)-\operatorname{Ad}(T_{yx})\lambda(x)$,
up to the fixed overall sign. At a mean-contour axial
representative $T_{yx}=I$. A representative that fixes
only a single comb need not have this property.
Order noncoarse sites by the hierarchy, putting a parent
before its descendants, and set the final coarse-node
values to zero. The resulting map on these site values
is triangular, with diagonal blocks
$\pm\operatorname{Ad}(T_{yx})$. These adjoint transports
have determinant one. The same is true of any transported
colour frames.
Consequently the actual axial coordinate $G_{\rm ax}$
satisfies
\begin{equation}\label{gi:eq-axial-constant}
 \sqrt{\det(R_0R_0^*)}\,
   |\det(G_{\rm ax}D|_{\ker R_0})|=c_{\eta,j},
\end{equation}
where $c_{\eta,j}>0$ depends on the meshes, hierarchy,
coordinate measures and fixed gauge weights, but not
on the background. The determinant one calculation is
in unweighted site coordinates; conversion by
\eqref{gi:eq-coordinate-measures} accounts for its
physical measure factors.

The linear coarea factor also equals the pivot-elimination
factor in a local fibre chart. To see this in orthonormal
coordinates, split $x=(u,p)$ so that
$Q=(Q_u,Q_p)$ with $Q_p$ invertible. The kernel is the
graph $p=-Q_p^{-1}Q_u u$. Its induced measure is
$\det(I+A^*A)^{1/2}du$, where $A=-Q_p^{-1}Q_u$.
Since
\[
 QQ^*=Q_p(I+AA^*)Q_p^*,\qquad
                  \det(I+AA^*)=\det(I+A^*A),
\]
the kernel measure divided by $\sqrt{\det QQ^*}$ is
$|\det Q_p|^{-1}du$. The determinant identity follows
by eliminating either diagonal block of
$\left(\begin{smallmatrix}I&-A^*\\A&I\end{smallmatrix}\right)$.
This is precisely the derivative of constraint
elimination. In the product Haar disintegration of
\cite[Proposition ``Local disintegration of product Haar
measure'']{BA}, the left exponential density at the
origin is field independent. Its central value and the
fixed coordinate scales multiply this linear integral
by constants; its nonconstant part is a nonquadratic
correction and remains outside the present calculation.

For clarity, these central Haar constants can be written
explicitly in the present colour convention. In a Pauli
basis of norm one, normalized $\operatorname{SU}(2)$ Haar
measure in $U=\exp(i x\cdot\sigma)$, $|x|<\pi$, is
\[
 \frac1{2\pi^2}\left(\frac{\sin|x|}{|x|}\right)^2d^3x.
\]
Indeed $\operatorname{SU}(2)$ is the unit three-sphere
$(\cos r,n\sin r)$. Multiplication by a unit quaternion
is orthogonal, so normalized spherical volume is Haar measure.
Its polar volume is
$\sin^2r\,dr\,d\Omega$, of total volume $2\pi^2$.
Euclidean polar measure is $r^2dr\,d\Omega$.
Thus the coordinate $U=e^{ih A}U_0$ has central Haar
density $(2\pi^2)^{-1}h^3d^3A$. The physical bond
Lebesgue measure of the pairing $h^4\langle A,A\rangle$
is $h^6d^3A$, giving the central ratio
\begin{equation}\label{gi:eq-central-haar-factor}
                         c_h=(2\pi^2)^{-1}h^{-3}.
\end{equation}
Before any gauge insertion, the linear fibre density
relative to prescribed target product Haar measure
therefore carries $\prod_{b\in E_f}c_{h_b}/
\prod_{c\in E_c}c_{h_c}$, multiplying the physical
delta integral. The sets $E_f,E_c$ here are all fine
and coarse link variables of that disintegration, not
only nonpivot variables. The pivot determinant is
already in the coarea factor. A subsequent rescaling
of the free integration coordinates contributes its
own Jacobian. These constants are retained for an
absolute normalizer and cancel only in a comparison
with the same coordinate and measure conventions.

\begin{proposition}[Normalized physical Gaussian]
\label{gi:prop-physical-normalizer}
Fix the spaces, meshes, coordinate measures and soft-gauge
normalization conventions. Suppose the following finite data
are given at each real background parameter $V$, including
a reference $V_*$.
The plain and straightened constraints satisfy
\[
 Q=Q_0+dS,\qquad R=R_0+SD,\qquad Q_0D=dR_0,
\]
and the hypotheses of the relative gauge identity,
with ineffective gauge directions treated by \cref{gi:cor-relative-gauge-stabilizer}, hold for both constraints. Their
physical quadratic form is represented by a common
symmetric $B$ with $BD=0$. A specified complete fibre
Hessian $B_{\rm fib}$ obeys
\begin{equation}\label{gi:eq-given-complete-hessian}
                     B=B_{\rm fib}-Q^*H_cQ
\end{equation}
for a symmetric coarse form $H_c$. Assume
\eqref{gi:eq-site-coisometry} and
\eqref{gi:eq-axial-constant}, and use the physical
coarea and central Haar conventions just described.
Let $\Psi(V)$ be the logarithm of the plain-fibre
Gaussian normalizer minus the same logarithm at $V_*$.
Then
\begin{align}
 \Psi(V)={}&
  -\tfrac12\log\left|\det\mathbb S
       (B_{\rm fib}+D(I-R^*R)D^*,Q)\right|
  +\log|\det\mathbb S(D^*D,R)|\nonumber\\
 &\quad-\left(\text{the same two terms at }V_*\right).
 \label{gi:eq-physical-normalizer}
\end{align}
In a nonlinear application the identification
$B_{\rm fib}=A''-\langle\mu,q''\rangle$ and the common
form in \eqref{gi:eq-given-complete-hessian} are supplied
data. This proposition does not prove their admission
or an extension to independent currents.
\end{proposition}

\begin{proof}
Write $c_{\rm Haar}$ for the product of the declared central
Haar factors, divided by the target Haar factors when that
is the target convention. Write $c_{\rm soft}>0$ for the
product of any prescribed fixed soft-gauge prefactors;
it is one when the quadratic gauge weights have no
prefactors. Both constants are independent of $V$.
The absolute plain-fibre normalizer
and the relative invariant satisfy separately
\[
 Z_{\rm plain}=c_{\rm Haar}c_{\rm soft}
                         I(B+G_{\rm ax}^*G_{\rm ax},Q_0),
 \qquad
 \mathcal J=c_{\eta,j}I(B+G_{\rm ax}^*G_{\rm ax},Q_0).
\]
Thus $Z_{\rm plain}=(c_{\rm Haar}c_{\rm soft}/c_{\eta,j})\mathcal J$.
Apply gauge and
constraint invariance, then use the background gauge on
the straightened side. Formula
\eqref{gi:eq-background-gauge-log} gives the two saddle
logarithms and the fixed power of $2\pi$. The row-penalty
invariance of the saddle determinant removes $Q^*H_cQ$
from the first quadratic form. Every removed constant
is independent of $V$, so subtraction at $V_*$ proves
\eqref{gi:eq-physical-normalizer}.
\end{proof}

The second saddle matrix stays invertible when $D$ has
global zero modes, provided it is injective on $\ker R$:
its compressed site form is then positive, and $R$ is
onto. Thus the formula uses finite constrained
determinants; an uncanceled whole-space site determinant
would not have the same property. For successive
integrations all measure constants must first be
retained in the exact composition formula. Cancellation
in a normalized logarithm uses the same reference and
the same coordinate conventions at the corresponding
stage.

\subsection{Scaling the action while keeping the gauge weight}

Let $V$ be a real $d$-dimensional Hilbert space and suppose
\[
 V=V_{\rm p}\oplus V_{\rm g},\qquad \dim V_{\rm g}=r.
\]
This is a direct sum, not necessarily orthogonal. Let
$K_A,K_G$ be symmetric quadratic forms such that $K_A$
vanishes on $V_{\rm g}$, $K_G$ vanishes on $V_{\rm p}$,
both have zero mixed terms, and each is positive on its
other summand. Thus $K=K_A+K_G>0$. These are the precise
splitting properties supplied by the complete fibre
Hessian and residual gauge coordinate in the linear
estimates; a partial Wilson Hessian need not have them.

\begin{proposition}[The two action-weight powers]
\label{gi:prop-action-gauge-powers}
For $\operatorname{Re}\tau>0$, put $K_\tau=\tau K_A+K_G$,
and continue powers from $\tau=1$. Then
\begin{align}
 \det K_\tau&=\tau^{d-r}\det K,\label{gi:eq-action-determinant}\\
 \int_Ve^{-u^{\mathsf T}K_\tau u/2}\,du
    &=\tau^{-(d-r)/2}\int_Ve^{-u^{\mathsf T}Ku/2}\,du,
       \label{gi:eq-action-integral}\\
 \int_V e^{(\tau-1)u^{\mathsf T}K_Gu/2}
                         d\mu_{(\tau K)^{-1}}(u)
    &=\tau^{r/2}.
       \label{gi:eq-gauge-compensation}
\end{align}
The last measure is the normalized Gaussian on all of $V$,
before any fluctuation test is inserted.
\end{proposition}

\begin{proof}
Choose fixed bases of the two summands. The two forms are
block diagonal in their concatenation, so multiplication
of the first block by $\tau$ multiplies the determinant
by $\tau^{d-r}$. The square of the basis-change Jacobian
is the same in numerator and denominator; orthogonality
is unnecessary. Integrating the two blocks for real
$\tau>0$ proves \eqref{gi:eq-action-integral}. For complex
$\tau$ the real part remains positive on compact subsets
of the right half-plane, so Gaussian domination and
holomorphic continuation prove the same identity there.
Finally the product of the exponential in
\eqref{gi:eq-gauge-compensation} and the density with
precision $\tau K$ has exponent $-u^{\mathsf T}K_\tau u/2$.
Its normalizing factor differs from that for $K$ by
$\tau^{d/2}$. Combining this with
\eqref{gi:eq-action-integral} gives $\tau^{r/2}$.
\end{proof}

Thus scaling the full Gaussian precision and compensating
the gauge weight contributes separately
$-d\log\tau/2$ and $r\log\tau/2$ to the logarithm of
the unnormalized integral. Their sum is
$-(d-r)\log\tau/2$. This identity contains no estimate
for a sharply truncated or interacting Gaussian integral.

\begin{proposition}[Successive residual gauge factors]
\label{gi:prop-gauge-stage-factors}
Let \(R_i:\mathcal H_{i-1}\to\mathcal H_i\), \(i=1,2\),
be onto maps of finite real site spaces with fixed Hilbert
measures. Let \(D_i:\mathcal H_i\to E_i\), \(i=0,1,2\),
be the gauge tangents of the three field spaces. Let
\(M:E_0\to E_1\) and \(N:E_1\to E_2\) be onto field rows
satisfying \(MD_0=D_1R_1\) and \(ND_1=D_2R_2\).
For \(i=1,2\) let \(Z_i\) be fixed real Hilbert target spaces
and let \(G_i:E_{i-1}\to Z_i\) be such that
\(G_iD_{i-1}|_{\ker R_i}:\ker R_i\to Z_i\) is invertible,
using the induced site measures on these kernels. Put
\(G_{\rm tot}=(G_1,G_2M):E_0\to Z_1\oplus Z_2\), with
the orthogonal product measure on the target. Then
\begin{equation}\label{gi:eq-gauge-stage-factors}
\begin{aligned}
 J_{R_2R_1}|\det(G_{\rm tot}D_0|_{\ker(R_2R_1)})|
  ={}& J_{R_1}|\det(G_1D_0|_{\ker R_1})|\\
     &{}\cdot J_{R_2}|\det(G_2D_1|_{\ker R_2})| .
\end{aligned}
\end{equation}
\end{proposition}

\begin{proof}
For a right inverse \(L_1\) of \(R_1\), the map
\(\Phi(u,v)=u+L_1v\) is an isomorphism
from \(\ker R_1\oplus\ker R_2\) to \(\ker(R_2R_1)\).
In the declared Hilbert measures the two-step coarea change
has factor \(J_{R_1}^{-1}J_{R_2}^{-1}\).
On the other hand first passing through the composite
fibre gives factor \(J_{R_2R_1}^{-1}|\det\Phi|\).
Thus
\[
                  |\det\Phi|=\frac{J_{R_2R_1}}{J_{R_1}J_{R_2}} .
\]
The lower block of \(G_{\rm tot}D_0\Phi(u,v)\) is
\(G_2D_1v\) by the first gauge-chain identity. Its upper
block is \(G_1D_0u+G_1D_0L_1v\). The map is triangular
with the two residual gauge maps on its diagonal. Taking
absolute determinants gives \eqref{gi:eq-gauge-stage-factors}.
\end{proof}

The number of residual site coordinates is additive across the
two stages. Thus any powers of \(2\pi\) used to normalize the
residual gauge Gaussians also multiply. Suppose the two
stages and the total form satisfy the hypotheses of
\cref{gi:thm-relative-gauge} or
\cref{gi:cor-relative-gauge-stabilizer}. If their gauge-fixed
precisions also obey the hierarchical splitting of
Section~\ref{gi:sec-successive}, multiplication of the two
relative gauge normalizers follows
from \eqref{gi:eq-successive-normalizers} and
\eqref{gi:eq-gauge-stage-factors}, with the stated measures
at the intermediate field and site spaces.

For a full periodic one-step $\operatorname{SU}(2)$ fibre with
$n$ coarse sites and blocking factor $L$, the admitted geometry
has $d=12(L^4-1)n$ free link coordinates and
$r=3(L^4-1)n$ residual site coordinates. There are $4L^4n$
stored positive fine bonds and $4n$ coarse bond constraints,
each with three colours; fixing the coarse nodes leaves
$(L^4-1)n$ site gauges. Thus $d-r=9(L^4-1)n$.
These counts apply to the full periodic fibre. A cut or
different gauge must use its actual dimensions and splitting.

\section{Complex Gaussian integrals}
\label{gi:sec-complex-gaussian}

Use orthonormal real coordinates on the selected fluctuation space. The
integrals below are over this fixed real coordinate space, with Lebesgue
measure \(du\). Matrix transposes are bilinear. Norms are Euclidean
operator norms, and real-part inequalities in a Gaussian exponent are
tested on real integration variables.

\begin{theorem}[The relative covariance sector]
\label{gi:thm-complex-sector}
Let \(D_0\) be a real positive definite matrix on \(\mathbb R^d\) with
\(\|D_0\|\le K\), where \(K>0\), and let \(D=D^{\mathsf T}\) satisfy
\[
 E=D_0^{-1/2}(D-D_0)D_0^{-1/2},\qquad r=\|E\|<\tfrac14 .
\]
Use the branch
\begin{equation}\label{gi:eq-complex-density}
 \log\det D=\log\det D_0+\operatorname{Tr}\log(I+E),\qquad
 d\mu_D(u)=(2\pi)^{-d/2}e^{-\frac12\log\det D}
                         e^{-u^{\mathsf T}D^{-1}u/2}\,du .
\end{equation}
Then \(\int d\mu_D=1\). For \(0\le\rho\le (8K)^{-1}\),
\begin{equation}\label{gi:eq-complex-growth}
 \int e^{\rho|u|^2}|d\mu_D(u)|
                 \le \exp\{C(r+K\rho)d\},
 \qquad
 |\log\det D-\log\det D_0|\le 2dr .
\end{equation}
Here \(C\) is absolute. The exponential transform for
\(\ell\in\mathbb C^d\) is
\begin{equation}\label{gi:eq-complex-exponential}
 \int e^{\ell^{\mathsf T}u}d\mu_D(u)
                         =e^{\ell^{\mathsf T}D\ell/2}.
\end{equation}
In particular all odd moments of linear forms vanish. For \(2m\) linear
forms their moment is the sum over pairings of their covariances
\(\ell_i^{\mathsf T}D\ell_j\), with absolute value at most
\((2m-1)!!\|D\|^m\prod_{i=1}^{2m}\|\ell_i\|\).
The absolute moment for \(j\ge1\) satisfies
\begin{equation}\label{gi:eq-complex-absolute-moment}
 \int \prod_{i=1}^j|\ell_i^{\mathsf T}u|\,|d\mu_D(u)|
       \le e^{Crd}(C\sqrt{Kj})^j\prod_{i=1}^j\|\ell_i\|.
\end{equation}
The measure on the zero-dimensional space has mass one.
\end{theorem}

\begin{proof}
The logarithm in \eqref{gi:eq-complex-density} is the norm-convergent
series for \(\log(I+E)\). Its trace gives the branch of the scalar
determinant logarithm because \(\det\exp B=\exp(\operatorname{Tr}B)\)
for a complex square matrix \(B\). The series bounds the trace
difference by \(d\sum_{k\ge1}r^k/k\le2dr\).

Make the real change \(u=D_0^{1/2}v\) and write
\[
 R_\rho=\operatorname{Re}(I+E)^{-1}-2\rho D_0.
\]
The real part here is entrywise. The matrix is real symmetric, and the
inverse series gives
\(\|R_\rho-I\|\le r/(1-r)+2\rho K\le 7/12\).
Thus \(R_\rho>0\). Real Gaussian integration gives the exact identity
\begin{equation}\label{gi:eq-absolute-sector}
 \int e^{\rho |u|^2}|d\mu_D(u)|
           =|\det(I+E)|^{-1/2}(\det R_\rho)^{-1/2}.
\end{equation}
The singular values of \(I+E\) are at least \(1-r\).
Every eigenvalue of \(R_\rho\) is at least \(1-r/(1-r)-2\rho K\).
Using \(-\log(1-x)\le Cx\) for \(0\le x\le 7/12\) proves the first
bound in \eqref{gi:eq-complex-growth}.

For fixed \(D_0\), the density and its exponential transform are
holomorphic in the independent complex symmetric entries of \(E\)
on the open norm ball. Indeed on each compact set there is an
integrable Gaussian majorant, also after multiplying by
\(e^{M|u|}\) for any fixed \(M\). The normalization and
\eqref{gi:eq-complex-exponential} hold on the real symmetric ball by
diagonalization and the real one-dimensional Gaussian integral.
For the exponential transform, the real identity first continues
in \(\ell\), using its entire dependence. The real symmetric ball
has a full set of real coordinates for the complex symmetric
space. The one-variable identity theorem in each coordinate on
a small product of discs, followed by connectedness of the norm
ball, proves both identities throughout the sector.

Differentiate \eqref{gi:eq-complex-exponential} in the scalar
coefficients of a finite list of linear forms. The exponential of
the quadratic form gives precisely the pairings. For the absolute
moment, the absolute density is its total mass times the real
Gaussian of covariance \(D_0^{1/2}R_0^{-1}D_0^{1/2}\), by
\eqref{gi:eq-absolute-sector} with \(\rho=0\). Its covariance norm
is at most \(3K/2\). A real linear functional of this Gaussian has
\(L^j\)-norm at most \(C\sqrt{Kj}\) times its coefficient norm.
This follows from the one-dimensional Gaussian even moments,
using the next even integer if necessary. Splitting a complex
linear functional into its real and imaginary parts gives the
same bound with a fixed factor. Apply H\"older's inequality with
exponent \(j\) to each of the \(j\) factors and use
\eqref{gi:eq-complex-growth}.
\end{proof}

\begin{corollary}[Parameters and localization cost]
\label{gi:cor-complex-parameters}
Let \(D(s)\) be holomorphic on a finite parameter domain and
satisfy the sector of \cref{gi:thm-complex-sector} about a fixed
real comparison matrix \(D_0\). Suppose \(F_s(u)\) is Borel in
the real variable \(u\), holomorphic in \(s\) for each \(u\), and
\(|F_s(u)|\le A e^{\rho |u|^2}\) on each closed parameter polydisc
under consideration, with \(0\le\rho\le(8K)^{-1}\).
A fixed bounded Borel test, including a characteristic function,
is allowed. Then \(I(s)=\int F_s\,d\mu_{D(s)}\) is holomorphic.
If the closed polydisc of radii \(t_\Delta>0\) centered at \(s\)
is in the given domain and \(r_*\) bounds \(\|E\|\) on it, then
\begin{equation}\label{gi:eq-parameter-derivative}
 |\partial_s^\nu I(s)|
 \le A\,\nu!\prod_{\Delta:\nu_\Delta>0}t_\Delta^{-\nu_\Delta}
                     \exp\{C(r_*+K\rho)d\}.
\end{equation}
In this estimate it suffices to vary the coordinates with
\(\nu_\Delta>0\). The same Cauchy bound applies to a finite set of
directional derivatives in the given field and current norms,
using only a polydisc of increments contained in the original
parameter domain. For the relative log determinant, the
same derivative bound uses \(2dr_*\) in place of the
expectation bound.

If each of \(N\) localization cubes carries at most \(aM^4\)
selected real coordinates, the exponential cost in
\eqref{gi:eq-parameter-derivative} is at most
\(\exp\{Ca(r_*+K\rho)M^4N\}\).
The corresponding determinant difference costs at most
\(2ar_*M^4N\) in its logarithm.
\end{corollary}

\begin{proof}
The lower bound on \(R_\rho\) in the proof of the theorem supplies
an integrable majorant uniform on the chosen polydisc.
Joint measurability follows by approximating the parameter by
finite grids and using its pointwise continuity.
Continuity and Morera's theorem, with Fubini on one coordinate
disc at a time, prove holomorphy. Cauchy's formula in the
differentiated variables proves \eqref{gi:eq-parameter-derivative}.
For the directional version compose with the finite affine map
from the coordinate discs to the given domain. For the determinant
apply the same formula to the trace-log bound in
\eqref{gi:eq-complex-growth}. The count is
\(d\le aM^4N\), with each selected coordinate counted once.
\end{proof}

On independent cube discs of radius \(\mathcal S>1\), the difference between \(s_\Delta=1\)
and \(s_\Delta=0\) costs at most \((R-1)^{-1}\) times the
supremum on \(|s_\Delta|=R\), for \(1<R<\mathcal S\).
Indeed the difference of the two Cauchy kernels is
\(1/[z(z-1)]\). Taking differences in distinct coordinates
multiplies these factors. This gives the same cost per
cube for mixed first differences of the Gaussian integral
or the relative logarithm.

In the matrix construction of \cite[Theorem F and Corollary 12.1]{MB},
\(D=D_W(s)\) is the covariance on the selected outer coordinates
\(W\), with \(\|E\|<1/4\) relative to its common real comparison.
The real coordinate count is the outer count on the wall and
held guards. The inner fine lattice and the full torus add no
dimension to this bound. The small sector budget there, chosen
after the localization scale, can in particular pay for the
factor \(M^4\) in the relative log determinant. The field and
current derivatives above are taken on contained discs in the
weighted domains of that construction. For independent cube
parameters of radius \(\mathcal S\), the Cauchy radii may be
chosen from the distance to its boundary.

At complete parameters, if \(A_Y\) is the precision on the
conditioned coordinate domain \(Y\), then
\(D_W=\chi_WA_Y^{-1}\chi_W^{\mathsf T}\). This is a marginal
after conditioning on \(Y\). For separated guarded components
the constructed covariance is block diagonal. Each block
inherits the relative sector: for its coordinate selector $J$,
the map $(JD_0J^{\mathsf T})^{-1/2}JD_0^{1/2}$ is a coisometry,
and the block relative perturbation is its compression of $E$.
The exponents split, and absolute convergence permits Fubini.
The full measure and each block measure have mass one, so their
normalizing determinant factors also multiply with the specified
branches. Thus the integrals of products of admissible tests factor.
These assertions use the prescribed covariance on
each declared domain. Independently chosen guarded extensions
need not have common marginals away from complete parameters.

\subsection{An auxiliary Gaussian and its local marginal}

\begin{proposition}[Eliminating an auxiliary Gaussian]
\label{gi:prop-auxiliary-gaussian}
Let \(C=C^{\mathsf T}\) and \(\Gamma\) be complex matrices of
sizes \(m\times m\) and \(m\times n\). Suppose \(C_0>0\) is real
symmetric, \(\Gamma_0\) is real, and
\[
 \|C_0\|+\|C_0^{-1}\|+\|\Gamma_0\|\le K,\qquad
 \|C-C_0\|+\|\Gamma-\Gamma_0\|\le\varepsilon .
\]
Put \(D=C+C\Gamma\Gamma^{\mathsf T}C\) and define \(D_0\) by the
same formula at the real reference. There are positive
\(\varepsilon_0,\rho_0,C_K\), depending only on \(K\), with the
following properties. If \(\varepsilon\le\varepsilon_0\) and
\(|F(u)|\le A e^{\rho|u|^2}\) for a real-variable Borel test
and \(0\le\rho\le\rho_0\), then the joint real-coordinate
integral below converges absolutely and satisfies
\begin{equation}\label{gi:eq-auxiliary-integral}
 \int d\mu_{I_n}(v)\,
 e^{-\frac12 v^{\mathsf T}\Gamma^{\mathsf T}C\Gamma v}
             \int e^{u^{\mathsf T}\Gamma v}F(u)\,d\mu_C(u)
       =\int F(u)\,d\mu_D(u).
\end{equation}
Its value at \(F=1\) is one, and
\begin{equation}\label{gi:eq-auxiliary-marginal-bound}
 \int |F(u)|\,|d\mu_D(u)|
                   \le A e^{C_K(\varepsilon+\rho)m}.
\end{equation}
The dimensions \(m,n\) do not enter the constants.
\end{proposition}

\begin{proof}
The joint precision on the real variables \((u,v)\) is
\[
 B=\begin{pmatrix}
    C^{-1}&-\Gamma\\
    -\Gamma^{\mathsf T}&I_n+\Gamma^{\mathsf T}C\Gamma
   \end{pmatrix}.
\]
At the reference its quadratic form is
\[
 (u-C_0\Gamma_0v)^{\mathsf T}C_0^{-1}(u-C_0\Gamma_0v)+|v|^2 .
\]
The triangular map \((u,v)\mapsto(u-C_0\Gamma_0v,v)\)
and its inverse have norm at most \(1+K^2\). Thus \(B_0\ge c_K I\).
The inverse identity for \(C\) and expansion of the remaining
products give \(\|B-B_0\|\le C_K\varepsilon\).
Choose the two radii so that this change and \(2\rho\) together
are less than \(c_K/2\). The real part of the exponent then
gives an integrable joint Gaussian majorant. Fubini applies to
\eqref{gi:eq-auxiliary-integral}.

Set \(A_v=I_n+\Gamma^{\mathsf T}C\Gamma\). It is in a fixed
relative precision sector about its real positive reference.
Its inverse is in the covariance sector of
\cref{gi:thm-complex-sector} after decreasing
\(\varepsilon_0\). Integrating \(v\) first, apply
\eqref{gi:eq-complex-exponential} with covariance \(A_v^{-1}\)
and force \(\Gamma^{\mathsf T}u\). The result is the factor
\[
 (\det A_v)^{-1/2}
       \exp\!\left(\tfrac12u^{\mathsf T}
             \Gamma A_v^{-1}\Gamma^{\mathsf T}u\right).
\]
There is no change in the real argument of \(F\). Direct
multiplication and the determinant obtained by eliminating either
off-diagonal block of
\(\left(\begin{smallmatrix} I_m&-C\Gamma\\
\Gamma^{\mathsf T}&I_n\end{smallmatrix}\right)\)
give
\[
 D^{-1}=C^{-1}-\Gamma A_v^{-1}\Gamma^{\mathsf T},
 \qquad \det D=\det C\,\det A_v .
\]
The determinant roots multiply at the real reference. Nonvanishing
and connectedness extend that equality to the small neighborhood.
The \(u\) density after integration is exactly \(d\mu_D\).
The convention \(n=0\) removes the auxiliary integral and factor.

Since \(D_0\ge C_0\), its inverse norm is at most \(K\);
also \(\|D_0\|\le C_K\).
Expanding \(D-D_0\) gives
\(\|D_0^{-1/2}(D-D_0)D_0^{-1/2}\|\le C_K\varepsilon\).
Thus \cref{gi:thm-complex-sector} gives normalization and
\eqref{gi:eq-auxiliary-marginal-bound}, after decreasing the
radii to fit that sector. The bound takes the absolute value
after forming the marginal, which accounts for the dimension
\(m\).
\end{proof}

Partition \(u\) into disjoint coordinate blocks \(u_b\).
For a set \(P\) of distinct blocks, thresholds \(B_b>0\), and
\(\gamma\ge0\) with \(\rho+\gamma/2\le\rho_0\),
\begin{equation}\label{gi:eq-auxiliary-failure}
 \int \prod_{b\in P}\mathbf1_{\{|u_b|\ge B_b\}}
              e^{\rho|u|^2}|d\mu_D(u)|
 \le
 \exp\!\left\{C_K(\varepsilon+\rho+\gamma/2)m
                   -\tfrac{\gamma}{2}\sum_{b\in P}B_b^2\right\}.
\end{equation}
On the event in the integrand \(|u|^2\ge\sum_{b\in P}B_b^2\).
Insert this inequality and apply
\eqref{gi:eq-auxiliary-marginal-bound} at growth
\(\rho+\gamma/2\). The argument needs no independence between
the tested coordinates. Fixed Borel restrictions and holomorphic
parameter dependence are covered by \cref{gi:cor-complex-parameters}.
For the rescaled pair
\((\tau^{-1}C,\tau^{1/2}\Gamma)\), the effective covariance is
\(\tau^{-1}D\). Use the positive square root at a real
\(\tau>0\) and its nearby continuation. On a fixed compact real
interval in \((0,\infty)\) the rescaled reference bounds are uniform,
and a sufficiently small complex displacement of \(\tau\) stays
in the same type of matrix neighborhood. When a region has
\(N\) localization cubes and \(m\le aM^4N\), its cost is
\(\exp\{C_Ka(\varepsilon+\rho)M^4N\}\). When the count is already
the number \(m\) of selected coordinates, that count is used
directly in \eqref{gi:eq-auxiliary-marginal-bound}.
A localization cube of side $M$ in coarse lattice units
contains at most $C(LM)^4$ fine bond-colour coordinates.
For this fine-coordinate auxiliary marginal one can
therefore take $a=CL^4$. This count is distinct from the
outer selected-coordinate count in the moving-wall
application above.

\section{Real Gaussian block restrictions}
\label{gi:sec-gaussian-cutoffs}

Let \(Z=(Z_b)_{b\in\mathcal I}\) be a centered real Gaussian
on finitely many Euclidean blocks, with coordinate covariance
\(\Gamma\ge0\). It can be represented as \(Z=A X\) with \(X\)
a standard Gaussian in \(\mathbb R^q\), allowing \(q=0\).
For \(B_b>0\), write
\[
 \mathcal B=\prod_b\{|z_b|\le B_b\},\qquad
 \nu=\mathcal L\bigl(Z\mid |Z_b|<B_b\text{ for every }b\bigr).
\]
This event has positive probability because its preimage contains
a neighborhood of zero. Let \(d_b=\operatorname{rank}\Gamma_{bb}\)
and \(\Gamma_{bb}\le\sigma_b^2 I\). If \(d_b>0\), take
\(\sigma_b>0\).

\begin{lemma}[Conditional supports and radial weights]
\label{gi:lem-gaussian-radial}
For \(d_b>0\), put \(V_b=\operatorname{ran}\Gamma_{bb}\).
Let $A_a$ be the block row of $A$ with output $Z_a$.
Let $W$ be an independent standard Gaussian on the possibly
zero-dimensional space $\ker A_b$, with its induced Euclidean
inner product, and set $A_a'=A_a|_{\ker A_b}$.
Let \(\Gamma_{bb}^{+}\) be the inverse on \(V_b\) extended
by zero.
Conditional on \(Z_b=z\in V_b\), the other blocks have a
version of their joint law given by
\begin{equation}\label{gi:eq-conditional-blocks}
 Z_a=\Gamma_{ab}\Gamma_{bb}^{+}z+A_a'W,\qquad a\ne b,
\end{equation}
For continuous even log-concave
functions \(0\le\psi_a\le1\) on the real blocks, the conditional
weight
\[
 w_b(z)=\mathbb E_W\prod_{a\ne b}
           \psi_a(\Gamma_{ab}\Gamma_{bb}^{+}z+A_a'W)
\]
is even and log-concave on \(V_b\). In particular
\(w_b(rn)\le w_b(sn)\) if \(r\ge s\ge0\) and \(|n|=1\) in \(V_b\).
\end{lemma}

\begin{proof}
Decompose \(X\) orthogonally into the row space of \(A_b\)
and \(\ker A_b\). These projections are independent standard
Gaussians on their subspaces. The projection onto the row space
is determined by \(z\) and equals
\(A_b^{\mathsf T}\Gamma_{bb}^{+}z\).
Applying \(A_a\) gives \eqref{gi:eq-conditional-blocks}.
No estimate of the regression matrix is needed.

We use a one-dimensional integration fact. Suppose
\(f,g,h\ge0\) are integrable continuous functions on the line,
positive on the interiors of their interval supports, and
\[
 h((1-t)x+ty)\ge f(x)^{1-t}g(y)^t,\qquad 0<t<1.
\]
If \(F=\int f\) and \(G=\int g\) are positive, let \(T\) be
the increasing quantile map from \(f/F\) to \(g/G\).
In the positive interior its derivative is
\(T'(x)=f(x)G/(F g(T(x)))\). Substitute
\((1-t)x+tT(x)\) in the integral of \(h\).
Its derivative is at least \(T'(x)^t\) by the arithmetic-geometric
mean inequality. Integrating first on compact intervals in the
positive interior and then exhausting it gives
\[
 \int h\ge\int f(x)^{1-t}g(T(x))^t T'(x)^t\,dx
                   =F^{1-t}G^t .
\]
If either integral vanishes, the same conclusion holds directly.
Applying this fact to slices of a jointly log-concave function
shows that integration in one coordinate preserves log-concavity.
The positive sets of these slices are intervals, and the slices
and their marginals are continuous by domination by a Gaussian
in the integrated variables. Iterating one coordinate at a time
therefore proves the marginal claim for the product of the
cutoffs and the Gaussian density of \(W\). That product is
log-concave by affine substitution and multiplication. It is
invariant under \((z,W)\mapsto(-z,-W)\). Hence \(w_b\) is even.
Writing \(sn\) as a convex combination of \(rn\) and \(-rn\)
gives the ray inequality, including when \(w_b(rn)=0\).
\end{proof}

For \(0<\delta\le1/2\), use smooth radial cutoffs
\(\psi_{B,\delta}\) that equal one on \(|z|\le(1-\delta)B\)
and vanish on \(|z|\ge B\). To construct them, put
\(\phi=0\) on \((-\infty,0]\), take
\[
 \phi(t)=\int_0^t(t-s)e^{-1/s^2}(1-s)^{-4}\,ds
       \quad(0<t<1),\qquad \phi=+\infty\text{ on }[1,\infty),
\]
and set \(\psi_{B,\delta}(z)=
\exp[-\phi((|z|/B-1+\delta)/\delta)]\).
The function \(\phi\) is nondecreasing and convex. It is flat
at zero and grows at least as \(c(1-t)^{-2}\) at one. Its
derivatives at one grow at most as inverse powers of \(1-t\),
so every derivative of the exponential vanishes there.
This proves smoothness of the cutoff at its two transition
spheres. Convex composition with the norm gives log-concavity.
As \(\delta\downarrow0\), the argument of \(\phi\) decreases at
every fixed \(|z|<B\), so the cutoffs increase to the strict-ball
indicator. Put
\[
 \chi_\delta=\prod_b\psi_{B_b,\delta}(Z_b),\qquad
 \nu_\delta[P]=\frac{\mathbb E[P(Z)\chi_\delta]}
                         {\mathbb E\chi_\delta}.
\]
The denominators are positive. For each polynomial, dominated
convergence gives \(\nu_\delta[P]\to\nu[P]\).

\begin{proposition}[Normalized cutoff flux]
\label{gi:prop-gaussian-cutoff-flux}
For each block define
\begin{equation}\label{gi:eq-smooth-flux}
 \beta_{b,\delta}=
 \frac{\mathbb E[|\nabla\psi_{B_b,\delta}(Z_b)|
                \prod_{a\ne b}\psi_{B_a,\delta}(Z_a)]}
      {\mathbb E\chi_\delta}.
\end{equation}
For \(d_b=0\) this is zero. If \(d_b>0\), then
\begin{equation}\label{gi:eq-flux-bound}
 \beta_{b,\delta}\le\beta_b^*=
 \frac{C_{d_b}}{B_b}
 \left[1+(B_b/\sigma_b)^{d_b}\right]
                      e^{-c_{d_b}B_b^2/\sigma_b^2}.
\end{equation}
Take \(\beta_b^*=0\) at rank zero. These bounds are independent
of the other blocks, of their number and dependencies, and of
any lower marginal eigenvalue. If \(B_b\ge\sigma_b>0\) the right
side is at most
\(2 C_{d_b}B_b^{d_b-1}\sigma_b^{-d_b}
 e^{-c_{d_b}B_b^2/\sigma_b^2}\).
\end{proposition}

\begin{proof}
A rank-zero block is identically zero and its cutoff is constant
there. For positive rank use \cref{gi:lem-gaussian-radial}
and polar coordinates \(z=rn\) in \(V_b\). Its intrinsic
Gaussian density in direction \(n\), omitting a common
determinant constant, is \(r^{d_b-1}e^{-q(n)r^2/2}\), where
\[
 q(n)=n^{\mathsf T}(\Gamma_{bb}|_{V_b})^{-1}n
                  \ge\sigma_b^{-2}.
\]
The radial gradient of the cutoff is also intrinsic in \(V_b\).
Set \(B=B_b\), \(a=(1-\delta)B\), and \(d=d_b\).
Let \(D_n\) denote the integral in \(r\) of this density times
\(\psi_{B,\delta}(rn)w_b(rn)\). On the plateau the cutoff is one.
For every \(r\in[a,B]\) ray monotonicity gives
\[
 D_n\ge w_b(rn)\,I(q(n),a),\qquad
 I(q,a)=\int_0^a s^{d-1}e^{-qs^2/2}\,ds .
\]
The negative radial derivative of the cutoff has integral one
on \([a,B]\). Thus the flux numerator in this direction is at
most \(D_n B^{d-1}e^{-q(n)a^2/2}/I(q(n),a)\).
This also holds if \(D_n=0\).

Integrating to \(\min(a,q^{-1/2})\) shows
\(I(q,a)\ge e^{-1/2}\min(a,q^{-1/2})^d/d\).
Since \(a\ge B/2\), the ratio is at most
\(C_d B^{-1}[1+(qB^2)^{d/2}]e^{-qB^2/8}\).
For \(qB^2\ge B^2/\sigma_b^2\), absorb the polynomial at
\(qB^2\) into half the exponential. This gives the bound in
\eqref{gi:eq-flux-bound}. Integration in \(n\) cancels the common
determinant constant and the integral of \(D_n\), which proves
the normalized estimate.
\end{proof}

The \(B_b^{-1}\) term matters at small radius. For an isotropic
\(d\)-dimensional block, the sharp radial boundary density divided
by the ball probability is asymptotic to \(d/B_b\) as
\(B_b/\sigma_b\to0\). This follows by substituting \(r=B_b t\)
in the two radial integrals.

\gineedspace{17\baselineskip}
\begin{corollary}[Integration by parts at a sharp block restriction]
\label{gi:cor-gaussian-cutoff-ibp}
For a linear statistic \(L=\sum_i\ell_i^{\mathsf T}Z_i\) and
polynomial \(P\), let
\begin{equation}\label{gi:eq-sharp-defect}
 \mathcal D_L(P)=\nu[LP]
               -\sum_{i,b}\ell_i^{\mathsf T}\Gamma_{ib}\nu[\nabla_bP].
\end{equation}
Then
\[
 |\mathcal D_L(P)|\le
 \|P\|_{\infty,\mathcal B}
                  \sum_{i,b}|\ell_i|\|\Gamma_{ib}\|\beta_b^* .
\]
For an empty block set all sums and the nonnegative budgets below are zero.
In particular, with positive amplitudes \(a_b\) and cell measures
\(m_b\), set
\[
 S=\sup_i\sum_bm_b\frac{\|\Gamma_{ib}\|}{a_i a_b},\qquad
 \varepsilon_\partial=\sup_b\frac{a_b\beta_b^*}{m_b}.
\]
The bound is at most
\(S\varepsilon_\partial\|P\|_{\infty,\mathcal B}
 \sum_i a_i|\ell_i|\).
\end{corollary}

\begin{proof}
Apply Gaussian integration by parts in \(X\) to \(P(AX)\)
times the smooth product cutoff. The two factors of \(A\) give
\(AA^{\mathsf T}=\Gamma\), regardless of its rank. The cutoff
derivative at block \(b\) contributes
\(\ell_i^{\mathsf T}\Gamma_{ib}\) times
\[
 \frac{\mathbb E[P(Z)\nabla\psi_{B_b,\delta}(Z_b)
                    \prod_{a\ne b}\psi_{B_a,\delta}(Z_a)]}
           {\mathbb E\chi_\delta}.
\]
Its norm is bounded by \(\|P\|_{\infty,\mathcal B}\beta_b^*\),
since the derivative and other cutoffs restrict to \(\mathcal B\).
The polynomial terms converge to the two terms of
\eqref{gi:eq-sharp-defect}. Thus the total defect has the
claimed bound, even if separate boundary contributions of
coincident constraints have no intrinsic sharp limit.
Insert \(a_i m_b/(a_i a_b)\) and the compensating factor
\(a_b/m_b\) in the displayed sum to obtain the weighted bound.
\end{proof}

Here \(\Gamma_{ib}\) is a coordinate covariance. If a positive
definite precision \(K\) is self-adjoint for the weighted pairing with
matrix \(M=\operatorname{diag}(m_b I)\), the coordinate density
uses precision \(MK\). When the inverse has physical kernel
\(\mathsf G\), the array identity is \(K^{-1}=\mathsf G M\),
so \((MK)^{-1}=\mathsf G\) is the coordinate covariance.
This explains the factor \(m_b^{-1}\) in the boundary budget.

\gineedspace{9\baselineskip}
\begin{corollary}[Linear statistics under block restrictions]
\label{gi:cor-gaussian-cutoff-moments}
For a real linear statistic \(L\) with unconditioned variance \(v^2\),
and any nonnegative increasing function \(f\) on \([0,\infty)\),
\[
                \nu[f(|L|)]\le \mathbb E f(v|N|),
               \qquad N\sim N(0,1).
\]
In particular \(\nu[L^{2m}]\le(2m-1)!!v^{2m}\),
\(\nu[e^{tL}]\le e^{t^2v^2/2}\) for real \(t\), and
\(\operatorname{Cov}_{\nu}(Z)\le\Gamma\).
\end{corollary}

\begin{proof}
If \(v=0\), both sides use \(L=0\). Otherwise condition \(X\)
on the single linear coordinate \(L=x\), using the orthogonal
decomposition in \cref{gi:lem-gaussian-radial}. With all
smooth cutoffs included, the resulting weight on the
one-dimensional Gaussian density is even and nonincreasing in
\(|x|\). Let \(R,R'\) be independent copies of \(v|N|\).
The product \((f(R)-f(R'))(w(R)-w(R'))\) is nonpositive.
For bounded \(f\), take expectations and divide by
\(\mathbb E w(R)>0\) to get the claimed inequality for
\(\nu_\delta\). Pass to \(\nu\) by dominated convergence and
then to unbounded \(f\) by truncation. The entire restricted
law is even. Apply the inequality to powers and to
\(\cosh(t|L|)\), and use the Gaussian moment identities.
Testing all real linear statistics proves the covariance order.
\end{proof}




\section{A holomorphic constrained coordinate chart}
\label{gi:sec-complex-chart}

A complex constraint does not define a real delta measure by
substitution into $\delta_F$. We instead continue a specified
real fibre chart. The normalized Gaussian theorems then act on
its fixed real coordinates. The chart Jacobian remains part of
the exact, unnormalized block integral.

Let $Q_0:E\to F$ be an onto real map, let $K_0=\ker Q_0$,
and let $L_0=Q_0^*(Q_0Q_0^*)^{-1}$. Write $i_0:K_0\to E$
for inclusion. Suppose $Q(z),T(z)$ depend holomorphically
on a parameter domain containing a real reference $z_0$,
with $Q(z_0)=Q_0$ and $T(z_0)$ real symmetric and positive
on $K_0$. Transposes below are bilinear continuations.
Assume $A(z)=Q(z)L_0$ is invertible and define
\begin{equation}\label{gi:eq-complex-fibre-chart}
 L(z)=L_0A(z)^{-1},\qquad
 Z(z)=(I-L(z)Q(z))i_0,\qquad
 \mathcal T(z)=Z(z)^{\mathsf T}T(z)Z(z).
\end{equation}
Then $QL=I$ and $QZ=0$. Relative to the decomposition
$E=i_0K_0\oplus L_0F$, the map $(u,b)\mapsto Zu+Lb$
has triangular blocks
$\left(\begin{smallmatrix}I&0\\-A^{-1}Qi_0&A^{-1}\end{smallmatrix}\right)$.
Its continued coarea density is therefore
\begin{equation}\label{gi:eq-complex-coarea-density}
             j(z)=J_{Q_0}^{-1}\det A(z)^{-1}.
\end{equation}
The orientation is chosen so that $j(z_0)>0$. Where $Q(z)$ is real
and $\det A(z)>0$, this is
the real measure factor in the $u$ coordinate. It is not
obtained by putting an absolute value into a complex determinant.

\begin{proposition}[A constrained integral in fixed coordinates]
\label{gi:prop-complex-chart}
In \eqref{gi:eq-complex-fibre-chart}, suppose
$T(z)^{\mathsf T}=T(z)$ and that
$\mathcal T(z)^{-1}$ lies in the relative covariance sector
of \cref{gi:thm-complex-sector} about
$\mathcal T(z_0)^{-1}$. For $b\in F_{\mathbb C}$ define
\[
 \mathcal I(z,b)=j(z)\int_{K_0}
   \exp\{-\tfrac12(Zu+Lb)^{\mathsf T}T(z)(Zu+Lb)\}\,du,
\]
where $u$ is real. Put
\begin{align}
 H(z)&=L-Z\mathcal T^{-1}Z^{\mathsf T}TL,
 &C(z)&=Z\mathcal T^{-1}Z^{\mathsf T},\nonumber\\
 \Sigma(z)&=L^{\mathsf T}TL
       -L^{\mathsf T}TZ\mathcal T^{-1}Z^{\mathsf T}TL.
 \label{gi:eq-complex-chart-matrices}
\end{align}
Then
\begin{equation}\label{gi:eq-complex-chart-integral}
 \mathcal I(z,b)=j(z)(2\pi)^{\dim K_0/2}
       \det(\mathcal T(z))^{-1/2}
                          e^{-b^{\mathsf T}\Sigma(z)b/2},
\end{equation}
with the positive-reference branch. The integral is
holomorphic in $z$ and entire in $b$. Its normalized
exponential test is
\begin{equation}\label{gi:eq-complex-chart-characteristic}
 \frac{j(z)}{\mathcal I(z,b)}
 \int_{K_0}e^{\ell^{\mathsf T}(Zu+Lb)}
          e^{-(Zu+Lb)^{\mathsf T}T(Zu+Lb)/2}\,du
       =e^{\ell^{\mathsf T}Hb+\ell^{\mathsf T}C\ell/2}.
\end{equation}
It gives all polynomial tests by differentiation in
$\ell\in E_{\mathbb C}$. In particular $QH=I$, $QC=0$,
and $Z^{\mathsf T}TH=0$.
\end{proposition}

\begin{proof}
The inverse covariance sector gives a strictly positive
real part of $u^{\mathsf T}\mathcal T u$, uniformly on
compact parameter sets. On compact sets of $b$ the cross
term is bounded by a constant times $|u|$. Gaussian
domination therefore proves absolute convergence and
holomorphic parameter dependence. Write
$g=Z^{\mathsf T}TLb$. The exponent is
$-u^{\mathsf T}\mathcal T u/2-g^{\mathsf T}u
-b^{\mathsf T}L^{\mathsf T}TLb/2$.
Apply \eqref{gi:eq-complex-exponential} to the fixed
real variable $u$ with covariance $\mathcal T^{-1}$ and
linear force $-g$. This proves
\eqref{gi:eq-complex-chart-integral}. With the additional
test, the force becomes $Z^{\mathsf T}\ell-g$ and the
constant term gains $\ell^{\mathsf T}Lb$. Expanding this
quadratic expression proves
\eqref{gi:eq-complex-chart-characteristic}. No integration
variable was shifted by a complex vector. The three
matrix identities follow directly from $QL=I$, $QZ=0$
and the definition of $H$.
\end{proof}

When $Q,T,b$ are real and $\det A>0$, this chart integral
is the original constrained integral. The untested
formula is independent of the chosen splitting on the common
connected parameter component containing the reference, when
the branches agree at the reference. Indeed another kernel
basis and right inverse have the form $Z'=ZM$ and
$L'=L+ZX$, with $M$ invertible and real at the reference.
Put $\epsilon_M=\operatorname{sgn}\det M(z_0)$.
The positive-reference coordinate density is multiplied by
$\epsilon_M\det M$, while the restricted quadratic form
becomes $M^{\mathsf T}\mathcal T M$. Its positive-reference
continued square root is multiplied by the same
$\epsilon_M\det M$. The
Schur form $\Sigma$ is unchanged by the shear $X$.
These algebraic identities prove agreement on that component.
On any other component of the overlap, the transition branches
must be matched separately; agreement at the reference alone
does not determine their relative sign there. No open set of
real points on a parameter curve is needed.
The exact determinant identity, including the chart density, is
\[
 j(z)^{-2}\det\mathcal T(z)
       =(-1)^{\dim F}\det\mathbb S(T(z),Q(z)).
\]
It follows by the same saddle congruence as
\eqref{gi:eq-saddle-det}, using $QZ=0$ and $QL=I$.
The right side is positive at the real reference.
The sign of a real saddle determinant must be accounted
for before selecting its positive-reference square root.

A fixed Borel test may be integrated on the real $u$
coordinates. At nonzero complex datum the normalized
density in this chart is the tilted measure
\begin{equation}\label{gi:eq-affine-tilt}
 e^{-g^{\mathsf T}u-g^{\mathsf T}\mathcal T^{-1}g/2}
                          d\mu_{\mathcal T^{-1}}(u),
 \qquad g=Z^{\mathsf T}TLb.
\end{equation}
The centered estimate cannot be used without this tilt.
For a test bounded by $F_0 e^{\rho|u|^2}$ with $F_0\ge0$ and
$0\le\rho\le(8K)^{-1}$, put $D=\mathcal T^{-1}$ and
$A_\rho=\operatorname{Re}D^{-1}-2\rho I$.
The sector proof gives $A_\rho>0$ and
$\|A_\rho^{-1}\|\le12K/5$. After taking the absolute
value, completion of a real square gives the additional
factor
\[
 \exp\!\left\{-\tfrac12\operatorname{Re}(g^{\mathsf T}Dg)
       +\tfrac12(\operatorname{Re}g)^{\mathsf T}
                     A_\rho^{-1}(\operatorname{Re}g)\right\}.
\]
It is at most $e^{CK|g|^2}$ because $\|D\|\le5K/4$.
Thus the normalized absolute integral is bounded by
$F_0\exp\{CK|g|^2+C(r+K\rho)\dim K_0\}$, including
the endpoint $\rho=(8K)^{-1}$.
Uniform bounds on a parameter polydisc give the Cauchy
derivative estimates in the same way. A Gaussian defined
instead on the centered chart $x=Hb+Zv$ has the centered
measure in its declared real $v$ coordinate. Arbitrary
Borel tests on these two complex charts are not identified.
They cannot be evaluated at complex-shifted points without
an extension and a proof of the resulting formula.
The finite determinant $j(z)$ is retained in the
unnormalized integral. A cost per localization cube for
a normalized selected Gaussian is the separate dimension
bound in \cref{gi:cor-complex-parameters}; the chart
identity alone gives no locality estimate for a global
coarea determinant.

\subsection{The physical normalizer in complex parameters}

\begin{corollary}[Complex physical parameters]
\label{gi:cor-physical-complex}
Suppose the data of \cref{gi:prop-physical-normalizer}
extend holomorphically to a connected physical parameter
chart containing an open neighborhood of $V_*$ in its
underlying real coordinate slice.
Require its real hypotheses on that neighborhood. Assume
the fibre Gaussians have the fixed-coordinate representations
and relative covariance sectors of
\cref{gi:prop-complex-chart}, with their transition branches
matched, and that the signed saddle determinants below
are nonzero and have holomorphic logarithms continued from
their positive real reference values. With bilinear
physical transposes, put
\begin{align*}
 s_{\rm f}(V)&=(-1)^{\dim F}\det\mathbb S
     (B_{\rm fib}+D(I-R^{\mathsf T}R)D^{\mathsf T},Q),\\
 s_{\rm s}(V)&=(-1)^{\dim\mathcal K}\det\mathbb S
                                      (D^{\mathsf T}D,R).
\end{align*}
Then the normalized physical Gaussian logarithm is
\[
 \Psi(V)=-\frac12\log\frac{s_{\rm f}(V)}{s_{\rm f}(V_*)}
                 +\log\frac{s_{\rm s}(V)}{s_{\rm s}(V_*)},
\]
with both logarithms zero at the reference. The parameters
are those of the physical chart; the conclusion asserts
no extension to independent currents off that chart.
\end{corollary}
\begin{proof}
The fixed-coordinate Gaussian formula and its domination
argument in \cref{gi:prop-complex-chart} prove holomorphy
of the normalized integral; that argument uses only the
stated covariance sector. Exponentiating the displayed
signed saddle expression gives a holomorphic function on
the chart. On a small connected real neighborhood of
$V_*$, $\det A$ is real, nonzero and has a fixed sign.
The continued chart integral differs from the positive
real integral by this fixed sign, which cancels on
normalization at $V_*$. Proposition~\ref{gi:prop-physical-normalizer}
therefore identifies the normalized chart integral with
the exponential of the displayed saddle expression.
An equality of holomorphic functions on an open real
coordinate neighborhood extends to a complex neighborhood:
all real partial derivatives of the difference vanish
there, and hence so do its complex Taylor coefficients.
The identity theorem on the connected chart proves their
equality. The displayed expression therefore supplies a
holomorphic logarithm of this normalized integral, equal
to zero at the reference. Nonvanishing of the integral
alone would not have guaranteed such a global logarithm.
\end{proof}

\gineedspace{8\baselineskip}
\section{The free action-density covariance}\label{gi:sec-free-density}

Use the Hermitian colour convention of
Section~\ref{gi:sec-real-gaussian} and write \(d_G=3\)
for its real dimension. The calculation below
uses \(d_G\) independent scalar colour components. Let \(g>0\)
be the coupling for which \(g^{-2}\) multiplies the Wilson
action in the exponent. On a raw link field \(A\), the
coboundary \(\mathrm d A\) is its unscaled oriented plaquette sum.
The quadratic action at the identity is
\[
                    \frac1{2g^2}\|\mathrm d A\|_2^2
\]
in counting coordinates. Indeed the logarithm of the
plaquette product of \(e^{itA_b}\) has first term
\(it(\mathrm d A)(p)\) and has trace zero. Thus the
\(t^2\) coefficient of \(1-\tfrac12\operatorname{Re}\operatorname{tr}U_p(t)\)
is \(\tfrac12|(\mathrm d A)(p)|^2\). This agrees with the
flat curl form of \cite[eqs.~(3.10)--(3.12), p.~392]{Bal85p}.

\subsection{The continuum curvature and its contractions}

Use the Fourier transform \(\widehat f(k)=\int e^{-ik\cdot x}f(x)\,dx\)
with inverse factor \((2\pi)^{-4}\). Define the potential as four
independent centred Gaussian test families per colour, with
covariance \(g^2 G\) in each family, where
\[
 \widehat G(k)=|k|^{-2},\qquad
 G(x)=\int_0^\infty(4\pi t)^{-2}e^{-|x|^2/(4t)}\,dt
                         =\frac1{4\pi^2|x|^2}\quad(x\ne0).
\]
The heat integral follows by Fourier inversion of the Gaussian
and the identity \(\int_0^\infty e^{-t|k|^2}\,dt=|k|^{-2}\).
Both identities hold against Schwartz tests by absolute
integration, using their decay and the local integrability
of \(|x|^{-2}\) and \(|k|^{-2}\) in four dimensions.
The covariance form is positive
because it is the \(L^2\) pairing of the Fourier tests multiplied
by \(|k|^{-1}\). In the separable real Hilbert completion for
this form, summing independent standard normals against the
orthonormal coordinates of a test constructs the family in
\(L^2\).

The curvature is \(F_{\mu\nu}=\partial_\mu A_\nu-\partial_\nu A_\mu\)
in the sense of test families. Write \(\mathsf K\) for its
unit-coupling covariance distribution in one colour and use
the orthonormal basis \(e_\mu\wedge e_\nu\), \(\mu<\nu\),
on \(\Lambda^2\mathbb R^4\). For \(P=(\mu,\nu)\) and
\(Q=(\rho,\sigma)\) the distribution is
\begin{equation}\label{gi:eq-curvature-kernel}
 \begin{split}
 \mathsf K_{PQ}={}&-\delta_{\nu\sigma}\partial_\mu\partial_\rho G
                   +\delta_{\nu\rho}\partial_\mu\partial_\sigma G\\
                &+\delta_{\mu\sigma}\partial_\nu\partial_\rho G
                   -\delta_{\mu\rho}\partial_\nu\partial_\sigma G.
 \end{split}
\end{equation}
One derivative acts in each covariance argument, which gives
the minus sign on the first term.

\begin{lemma}[The curvature tensor away from coincidence]
\label{gi:lem-curvature-reflection}
For \(u\ne0\), put \(\widehat u=u/|u|\) and
\(R_u=I-2\widehat u\widehat u^{\mathsf T}\). Then
\[
       \mathsf K(u)=\pi^{-2}|u|^{-4}\Lambda^2R_u,\qquad
       \|\mathsf K(u)\|_{\mathrm{HS}}^2=6\pi^{-4}|u|^{-8}.
\]
Here \(\|\cdot\|_{\mathrm{HS}}\) is the Hilbert--Schmidt norm
in the six-dimensional two-form basis.
\end{lemma}

\begin{proof}
The Hessian of \(G\) is
\((4\widehat u\widehat u^{\mathsf T}-I)/(2\pi^2|u|^4)\).
For a symmetric matrix \(M\), its action on each of the
two factors of a wedge, added together, sends
\(v\wedge w\) to \(Mv\wedge w+v\wedge Mw\).
Equation~\ref{gi:eq-curvature-kernel} is the negative
of this action for \(M=\nabla^2G\). Since
\(\widehat u\widehat u^{\mathsf T}\) has rank one, the
product term with this matrix acting on both factors
vanishes. Expanding \(\Lambda^2R_u\) therefore gives
the stated identity. The reflection is orthogonal, as is
its two-form action.
\end{proof}

To define the separated density correlations, fix an even,
nonnegative smooth function $\chi$ of integral one supported
in the unit ball and set
$\chi_\varepsilon(x)=\varepsilon^{-4}\chi(x/\varepsilon)$.
Let \(F_\varepsilon(x)\) be the curvature
tested against its translate to \(x\). Use the centred density
\[
 q_\varepsilon(x)=\frac12\sum_{\alpha=1}^{d_G}\sum_P
  \bigl(F_{\varepsilon,P}^{\alpha}(x)^2
          -\mathbb E[F_{\varepsilon,P}^{\alpha}(x)^2]\bigr).
\]
For each \(\varepsilon>0\), integrate this field against
continuous compactly supported tests by \(L^p\) Riemann sums,
\(1\le p<\infty\). The smoothed covariance is continuous.
Gaussian moment bounds then give \(L^{2p}\) continuity of
each curvature component and \(L^p\) continuity of its square.

\begin{proposition}[Separated connected Gaussian densities]
\label{gi:prop-free-continuum-cycles}
For \(n\ge2\), the mixed cumulant of the regularized density
has a limit on every collection of pairwise separated compact
supports. For noncoincident points \(x_1,\ldots,x_n\) its kernel is
\begin{equation}\label{gi:eq-free-cycle}
 C_n(x_1,\ldots,x_n)=\frac{d_Gg^{2n}}{2n}
  \sum_{\sigma\in S_n}
  \operatorname{tr}\!\left(
    \prod_{j=1}^n\mathsf K(x_{\sigma(j)}-x_{\sigma(j+1)})\right),
 \qquad \sigma(n+1)=\sigma(1).
\end{equation}
The product follows the displayed order. In particular the
two-point smearing limit is \(C_Gg^4\mathcal K(f,h)\), where
\begin{equation}\label{gi:eq-free-coefficient}
       C_G=\frac{3d_G}{\pi^4}=\frac9{\pi^4},\qquad
       \mathcal K(f,h)=\iint f(x)h(y)|x-y|^{-8}\,dx\,dy.
\end{equation}
Every odd kernel \(C_n\), \(n\ge3\), vanishes on the separated
configuration space. At four distinct collinear points,
\begin{equation}\label{gi:eq-free-collinear-four}
 C_4(x_1,x_2,x_3,x_4)=\frac{6d_Gg^8}{\pi^8}
       \sum_{\text{three unoriented cycles}}
          \prod_{\{i,j\}\,\text{in the cycle}}|x_i-x_j|^{-4}.
\end{equation}
This is strictly positive and remains positive on the product
of sufficiently small disjoint neighborhoods of the four points.
\end{proposition}

\begin{proof}
For any real finite Gaussian vector with covariance \(C\),
write it as \(C^{1/2}\) times a standard Gaussian. If \(T_i\)
are real symmetric matrices, the logarithmic moment generating
function of its centred half-quadratic forms is
\[
 -\frac12\log\det\bigl(I-C\textstyle\sum_i z_iT_i\bigr)
                 -\frac12\sum_i z_i\operatorname{tr}(CT_i).
\]
Diagonalizing the symmetric quadratic form for small real \(z\)
and integrating each coordinate proves the formula. Sylvester's
determinant identity also covers singular \(C\). The resulting
power series at zero is obtained by integrating the convergent
geometric series for the inverse in the derivative of the
logarithm. For \(n\ge2\), the mixed derivative of order \(n\)
comes from its term \(\operatorname{tr}((C\sum_i z_iT_i)^n)/(2n)\).
For curvature at \(n\) points use the coordinate projections for
the \(T_i\). Each permutation gives one ordered cycle of
covariance blocks. Independence multiplies the expression
by \(d_G\).

For separated supports the smoothed cross-covariance blocks
converge uniformly, with every fixed derivative, to
\(g^2\mathsf K\). Indeed their convolutions sample a fixed
compact neighborhood away from coincidence, on which each
derivative of the kernel is uniformly continuous. Differentiation
under the convolution and the unit integral of the mollifier
give the asserted convergence. Integrating the finite cycle formula and taking
this limit proves \eqref{gi:eq-free-cycle}. For \(n=2\) the
two permutations give the squared Hilbert--Schmidt norm with
prefactor \(d_Gg^4/2\). Lemma~\ref{gi:lem-curvature-reflection}
gives \eqref{gi:eq-free-coefficient}.

On two-forms the Hodge star \(J\) satisfies \(J^2=I\) and
\(J\Lambda^2R=(\det R)\Lambda^2R\,J\) for \(R\) orthogonal.
This follows by applying \(R\) to the defining wedge identity
for the star and tracking its action on the oriented volume.
Each \(\mathsf K(u)\) thus anticommutes with \(J\), so
conjugation by \(J\) negates an odd product and preserves its trace.
At collinear points all the reflections agree, since
\(R_u=R_{-u}\). The product of four two-form actions is the
identity, whose trace is six. Each unoriented four-cycle has
eight representatives in \(S_4\), giving
\eqref{gi:eq-free-collinear-four}. Smoothness at separated
arguments then gives positivity on the stated neighborhoods.
\end{proof}

The separated limits concern mixed cumulants. In fact for a
nonnegative nonzero test \(f\), the variance of
\(\int f(x)q_\varepsilon(x)\,dx\) grows at least as
\(c_f g^4\varepsilon^{-4}\). Choose a ball on which \(f\)
has a positive lower bound. In its interior, and for
\(M\varepsilon\le |x-y|\) with \(M\) sufficiently large,
the convolution error in \(\mathsf K(x-y)\) is at most
\(C\varepsilon|x-y|^{-5}\), by its first derivative.
Its squared Hilbert--Schmidt norm is thus bounded below by
\(c|x-y|^{-8}\). The Wick variance integrates this
bound on a shell with fixed outer radius inside the ball.
Its radial integral is a positive multiple of
\(\int_{M\varepsilon}^{r_0}r^{-5}\,dr\).

\subsection{The lattice projection and its physical kernel}

Let \(a>0\), \(\ell=ma\) with \(m\) a positive integer, and
\(\Lambda_{a,\ell}=a\mathbb Z^4/\ell\mathbb Z^4\).
The coordinates are based positive bonds and plaquettes.
Write \(p=(x,P)\), \(P=(\mu,\nu)\) with \(\mu<\nu\), and
\(x_p=x+ac_P\), where \(c_P=(e_\mu+e_\nu)/2\).
Use the exact-form projection \(P_{\rm ex}\) of
Section~\ref{gi:sec-free-slices} in one colour, also on lattices
with no blocking step. On \((\ker\mathrm d)^\perp\) the map
\(\mathrm d\) is a linear isomorphism onto \(\mathcal E\).
It sends the normalized Gaussian for the raw quadratic
action to the Gaussian with covariance \(g^2P_{\rm ex}\).
The change of coordinates has a constant Jacobian, which
cancels in this normalized law.

For independent raw colour curvatures \(F_p^\alpha\) with
this covariance, their physical values are \(a^{-2}F_p^\alpha\).
The physical unit-coupling covariance entry is consequently
\begin{equation}\label{gi:eq-lattice-curvature-kernel}
            \mathsf K_{a,\ell}(p,q)=a^{-4}P_{\rm ex}(p,q).
\end{equation}
The centred quadratic observable for a real test \(f\) on the
torus is
\begin{equation}\label{gi:eq-free-lattice-density}
 Q_a(f)=\frac12\sum_{\alpha=1}^{d_G}\sum_p f(x_p)
                \bigl((F_p^\alpha)^2-\mathbb E[(F_p^\alpha)^2]\bigr).
\end{equation}
In physical coordinates the \(a^4\) integration weight cancels
the \(a^{-4}\) in the square of the physical curvature.
The finite Gaussian calculation in
Proposition~\ref{gi:prop-free-continuum-cycles}, with the \(T_i\)
given by diagonal test multipliers on plaquettes, yields
\begin{equation}\label{gi:eq-free-lattice-wick}
 \operatorname{Cov}(Q_a(f),Q_a(h))
     =\frac{d_Gg^4}{2}\,a^8\sum_{p,q}
                      f(x_p)h(x_q)\mathsf K_{a,\ell}(p,q)^2 .
\end{equation}
For a mixed cumulant of order \(n\), the same calculation gives
a sum over the permutations in \eqref{gi:eq-free-cycle}
and all \(p_1,\ldots,p_n\). Its summand has
\(a^{4n}\prod_i f_i(x_{p_i})\) times the product of the entries
\(\mathsf K_{a,\ell}(p_{\sigma(j)},p_{\sigma(j+1)})\).
The prefactor is \(d_Gg^{2n}/(2n)\). These are real scalar
entries, including when the plane labels differ.

Let \(\delta=2\pi/\ell\) and choose the finite frequency set
\(\mathcal B_a=\delta\mathbb Z^4\cap[-\pi/a,\pi/a)^4\).
For \(k\in\mathcal B_a\) put
\[
 t_{a,\mu}(k)=a^{-1}(e^{iak_\mu}-1),\qquad
 \omega_a(k)=|t_a(k)|^2,\qquad B_a(k)v=t_a(k)\wedge v.
\]
The symbol \(B_a\) represents the difference \(a^{-1}\mathrm d\),
which has the same range as \(\mathrm d\).
The wedge norm identity gives
\(B_a^\dagger B_a=\omega_a I-t_at_a^\dagger\).
Here \(\dagger\) is the Hermitian adjoint of a Fourier matrix.
As \(B_a t_a=0\), the projection symbol is
\begin{equation}\label{gi:eq-curvature-projection}
    \Pi_a(k)=B_a(k)B_a(k)^\dagger/\omega_a(k)\quad(k\ne0),
                    \qquad \Pi_a(0)=0 .
\end{equation}
The rank is three at each nonzero frequency. A constant
two-form has zero exact component, including on the one-site
torus. The unitary discrete Fourier transform has factor
\(m^{-2}\) in both directions. Inserting it in
\eqref{gi:eq-lattice-curvature-kernel} gives, for \(p=(x,P)\)
and \(q=(y,Q)\),
\begin{equation}\label{gi:eq-lattice-curvature-fourier}
  \mathsf K_{a,\ell}(p,q)
  =\left[\ell^{-4}\sum_{k\in\mathcal B_a}
                   e^{ik\cdot(x-y)}\Pi_a(k)\right]_{PQ}.
\end{equation}
In particular the finite Fourier coefficient is
\(a^{-4}m^{-4}=\ell^{-4}\).

The continuum multiplier \(\Pi_{\mathrm{cont}}(k)\) is
\eqref{gi:eq-curvature-projection} with \(t_0=ik\) and
\(\omega_0=|k|^2\), and \(\Pi_{\mathrm{cont}}(0)=0\) on the
reciprocal grid.
Its inverse transform on \(\mathbb R^4\) is
\eqref{gi:eq-curvature-kernel}.
On the torus write \(\mathsf K_\ell\) for the distribution with
Fourier series
\(\ell^{-4}\sum_{k\in\delta\mathbb Z^4}
e^{ik\cdot z}\Pi_{\mathrm{cont}}(k)\).
The zero-frequency convention matters even away from
coincidence. The traces of the raw lattice matrix in
\eqref{gi:eq-lattice-curvature-fourier} and of \(\mathsf K_\ell\)
are, respectively,
\[
  3\bigl(a^{-4}\mathbf 1_{\{z=0\}}-\ell^{-4}\bigr),
                 \qquad 3\bigl(\delta_{\rm per}(z)-\ell^{-4}\bigr).
\]
These identities follow by taking the trace of the projection
symbol. Here \(z=x-y\) and \(\delta_{\rm per}\) is the unit point
mass at the origin of the torus.

\begin{proposition}[Separated comparison of curvature kernels]
\label{gi:prop-free-lattice-kernel}
Let \(r\) be the Euclidean distance on the torus from \(x_p\)
to \(x_q\). For \(r\ge10a\), if \(u\) is the short displacement
from \(x_q\) to \(x_p\), then
\begin{equation}\label{gi:eq-free-fixed-torus-error}
   \left|\mathsf K_{a,\ell}(p,q)-\mathsf K_{\ell,PQ}(u)\right|
             \le C r^{-4}(a/r)^{1/2}.
\end{equation}
For \(10a\le r\le\ell/10\) one has
\begin{equation}\label{gi:eq-free-volume-mesh-error}
   \left|\mathsf K_{a,\ell}(p,q)-\mathsf K_{PQ}(u)\right|
      \le C r^{-4}\bigl[(a/r)^{1/2}+(r/\ell)^4\bigr].
\end{equation}
The constant is independent of the orientation labels, the
mesh and the period. The smooth torus kernel in
\eqref{gi:eq-free-fixed-torus-error} is periodic in \(u\).
\end{proposition}

\begin{proof}
We first obtain the continuum torus estimates. Its heat kernel is
the sum of translates by \(\ell\mathbb Z^4\) of
\(p_t(x)=(4\pi t)^{-2}e^{-|x|^2/(4t)}\), since integrating
this sum against each Fourier mode gives
\(\ell^{-4}e^{-t|k|^2}\). Integrating the heat kernel with
its mean \(\ell^{-4}\) subtracted gives the periodic Green
distribution \(G_\ell\) of zero mean. Inserting \(G_\ell\)
into \eqref{gi:eq-curvature-kernel} produces \(\mathsf K_\ell\).

At period one and \(|z|\le1/4\), the difference \(G_1-G\)
is smooth. On \(0<t\le1\) all image terms other than the
central one, with any fixed number of derivatives, have an
integrable bound \(C t^{-M}e^{-c/t}\) for suitable \(M\).
This follows from the lower bound \(|z+m|\ge c|m|\) on each
noncentral image. On \(t\ge1\), the nonzero torus Fourier
modes decay exponentially with all fixed derivatives, and
the corresponding free heat derivatives are bounded by
\(C_\alpha t^{-2-|\alpha|/2}\). The mean subtraction over
\((0,1]\) is constant. Thus the difference has bounded
derivatives of all fixed orders on the ball. On the rest of the
unit torus away from zero the same split gives smoothness
and bounded derivatives on closed sets. Scaling
\(G_\ell(z)=\ell^{-2}G_1(z/\ell)\) shows, with
\(r_z=\operatorname{dist}(z,\ell\mathbb Z^4)>0\), that
\begin{equation}\label{gi:eq-free-torus-green-bounds}
 \begin{aligned}
 \|\partial^\alpha\mathsf K_\ell(z)\|
       &\le C_\alpha r_z^{-4-|\alpha|} &&(r_z>0),\\
 \|\partial^\alpha(\mathsf K_\ell-\mathsf K)(z)\|
       &\le C_\alpha\ell^{-4-|\alpha|}
                         &&(0<|z|\le\ell/4).
 \end{aligned}
\end{equation}

We compare the Fourier sums at the raw displacement \(z=x-y\).
The bound \(|t_a(k)|\asymp|k|\) on the Brillouin cube follows
from the sine formula for \(|t_a|\).
Differentiating \(t_a\) gives
\(|\partial^\alpha t_a|\le C_\alpha|k|^{1-|\alpha|}\)
away from zero, since \(a|k|\) is bounded on that cube.
The product and quotient rules in \eqref{gi:eq-curvature-projection}
therefore give
\(\|\partial^\alpha\Pi_a(k)\|\le C_\alpha|k|^{-|\alpha|}\).
The same estimates hold for \(\Pi_{\mathrm{cont}}\). Taylor
expansion of \(t_a\) yields
\(\|\Pi_a(k)-\Pi_{\mathrm{cont}}(k)\|\le Ca|k|\) on
\(a|k|\le c_0\), for a fixed small numerical \(c_0\).

Choose a smooth frequency cutoff equal to one for \(|k|\le R\)
and zero for \(|k|\ge2R\). Assume for now that \(R\ge18\delta\)
and \(2aR\le c_0\). On its support the two reciprocal grids
coincide. Counting disjoint cubes of side \(\delta\) around
the grid points shows that the low-frequency difference has
norm at most \(CaR^5\), with the Fourier factor \(\ell^{-4}\).
For the high-frequency lattice symbol, apply the ninth
forward frequency difference in a coordinate \(\mu\) with
\(|z_\mu|\ge r_z/2\) in a short representative. The cutoff
extends periodically on the Brillouin torus. The derivative
estimates above and the product rule bound this difference at
\(k\) by
\[
   C\delta^9\bigl(R+\operatorname{dist}(k,(2\pi/a)\mathbb Z^4)\bigr)^{-9}.
\]
Indeed the difference is the integral of the ninth derivative
over nine shifts of length at most \(\delta\), and \(R\)
dominates the sum of those lengths. Summation by parts has
factor \((e^{-i\delta z_\mu}-1)^9\) on the Fourier sum.
The raw \(z/a\) is integral, so the sum has periodic phase
on this frequency torus. The modulus of the factor is at
least \((2r_z/\ell)^9\), by the sine bound for
\(|z_\mu|\le\ell/2\). The ninth difference has absolute sum
at most \(C\delta^9R^{-5}\) after multiplying by \(\ell^{-4}\).
This follows by comparing the grid sum of
\((R+|k|)^{-9}\) with its integral. The same argument applies
to the high-frequency continuum symbol on the full reciprocal
grid. For its distributional Fourier series the multiplication
by \((e^{-i\delta z_\mu}-1)^9\) gives an absolutely convergent
series of ninth differences. Dividing away from the zero set
of this factor proves the pointwise bound there.
Both high-frequency terms are thus at most \(C r_z^{-9}R^{-5}\).

Take \(R=D r_z^{-1}(r_z/a)^{1/10}\) with a fixed \(D\)
large enough that \(R\ge18\delta\), since \(r_z\le\ell\).
For \(r_z/a\) above a fixed numerical threshold the support
condition \(2aR\le c_0\) holds. The sum of the two errors is
bounded by \(C r_z^{-4}(a/r_z)^{1/2}\).
For the remaining nonzero raw displacements, the lattice
Fourier sum has norm at most \(a^{-4}\) and
\eqref{gi:eq-free-torus-green-bounds} bounds the torus
kernel by \(Cr_z^{-4}\). Here \(1\le r_z/a\) is bounded
above by the fixed threshold. Increasing \(C\) gives the
same claimed raw-displacement bound.

The shift \(a(c_P-c_Q)\) from origins to centres has length
at most \(a\). When \(r\ge10a\), the shifted segment has
torus distance at least \(r-a\) from zero. The first derivative
bound in \eqref{gi:eq-free-torus-green-bounds} gives an error
of at most \(Ca r^{-5}\) for that shift.
This proves \eqref{gi:eq-free-fixed-torus-error}. Applying the
second estimate in \eqref{gi:eq-free-torus-green-bounds}
at \(u\) proves \eqref{gi:eq-free-volume-mesh-error}.
\end{proof}

\subsection{Separated tests and a coarsened Gaussian}

\begin{proposition}[Limits of the separated lattice cumulants]
\label{gi:prop-free-density-limits}
For \(n\ge2\) and continuous tests on the torus with pairwise
separated supports, the mixed cumulant of \(Q_a(f_i)\) converges,
as \(a\downarrow0\) through subdivisions of \(\ell\), to the
integral of \eqref{gi:eq-free-cycle} with \(\mathsf K_\ell\)
in place of \(\mathsf K\).
For periodizations of tests in a fixed compact subset of
\(\mathbb R^4\) with positive pairwise separation, it
converges to the integral of
\eqref{gi:eq-free-cycle} along every sequence with
\(a\downarrow0\), \(\ell\to\infty\) and \(\ell/a\) integral.
The same Euclidean integral is obtained by taking the
fixed-volume limit first and then increasing \(\ell\).

For a quantitative version, let \(n\) be fixed and let the
tests be real Lipschitz functions, each supported in a ball
of radius \(d\) in a common Euclidean lift. Suppose every
distance between a point of one support and a point of a
distinct support lies in \([d,C_{\rm sh}d]\),
\(10a\le d\) and \(C_{\rm sh}d\le\ell/10\), and that
\(\operatorname{Lip}(f_i)\le M\|f_i\|_\infty/d\).
The difference between the lattice cumulant and the integral
of \eqref{gi:eq-free-cycle} is bounded by
\begin{equation}\label{gi:eq-free-higher-error}
 C_{n,M,C_{\rm sh}}d_Gg^{2n}\prod_i\|f_i\|_\infty
                     \bigl[(a/d)^{1/2}+(d/\ell)^4\bigr].
\end{equation}
For odd \(n\ge3\) this is an absolute bound on the lattice
cumulant.
\end{proposition}

\begin{proof}
Use the finite cycle sum following \eqref{gi:eq-free-lattice-wick}.
Each edge of a contributing cycle joins distinct supports.
The convergence of every kernel entry is uniform there by
Proposition~\ref{gi:prop-free-lattice-kernel}.
The plaquette counts in the sum, after multiplication by
\(a^{4n}\), are bounded on a fixed compact support stock.
For each set of plane labels the limit is the Riemann integral
on the corresponding shifted grids. The integrand is continuous
after being assigned the value zero in a neighborhood of each
partial diagonal, since the product of the tests vanishes
there. Summing the plane labels gives the trace in the cycle.
The fixed-volume and joint limits follow from the two kernel
estimates. The periodic-to-Euclidean difference after taking
the fixed-volume limit tends to zero by
\eqref{gi:eq-free-torus-green-bounds}.

For the rate, write the difference of the cycle products as
a sum in which one kernel changes in each term. All kernels
are \(O(d^{-4})\) on the contributing configurations.
By \eqref{gi:eq-free-volume-mesh-error} the changed kernel
has the additional factor \(C[(a/d)^{1/2}+(C_{\rm sh}d/\ell)^4]\).
The \(a^{4n}\)-weighted
count of sites on all test supports is \(O(d^{4n})\).
This bounds the lattice kernel error by the right side of
\eqref{gi:eq-free-higher-error}. The Riemann error of the
Euclidean integrand has the additional factor \(C_{n,M}a/d\)
in place of the bracket. Indeed each kernel has first
derivative \(O(d^{-5})\) on the cells of the contributing
stock, enlarged by one mesh, and the same product estimate
uses the stated Lipschitz bounds of the tests. The total
volume of these cells is \(O(d^{4n})\).
Since \(a/d\le1\), this last error is absorbed in the
displayed bound. The odd case uses the exact continuum
cancellation in Proposition~\ref{gi:prop-free-continuum-cycles}.
\end{proof}

\gineedspace{11\baselineskip}
\begin{corollary}[Finite volume for continuous separated tests]
\label{gi:cor-continuous-volume}
Let $f,h$ be nonzero nonnegative continuous compactly supported
tests in a common Euclidean lift. Suppose all cross-distances
between their supports lie in $[d,C_{\rm sh}d]$, where
$d>0$, $C_{\rm sh}\ge1$, and
$C_{\rm sh}d\le\ell/4$. Denote the limiting torus two-point
covariance of their periodizations, as in
\cref{gi:prop-free-density-limits}, by
$\mathcal C_\ell(f,h)$. Then
\[
 \left|\frac{\mathcal C_\ell(f,h)}
 {C_Gg^4\mathcal K(f,h)}-1\right|
       \le C_{C_{\rm sh}}(d/\ell)^4.
\]
No Lipschitz or lower-mass assumption is required for this
finite-volume comparison of the limiting kernels.
\end{corollary}
\begin{proof}
The two-point Wick formula gives
\[
 \mathcal C_\ell(f,h)=\frac{d_Gg^4}{2}
       \iint f(x)h(y)\|\mathsf K_\ell(x-y)\|_{\rm HS}^2\,dx\,dy.
\]
Equation~\ref{gi:eq-free-torus-green-bounds} and
$\|\mathsf K(z)\|_{\rm HS}\le C|z|^{-4}$ give
\[
 \left|\|\mathsf K_\ell(z)\|_{\rm HS}^2
              -\|\mathsf K(z)\|_{\rm HS}^2\right|
                \le C\ell^{-4}|z|^{-4}
       \quad(0<|z|\le\ell/4).
\]
Multiply by $f(x)h(y)\ge0$ and use
$|x-y|^{-4}\le(C_{\rm sh}d)^4|x-y|^{-8}$.
The denominator is positive because both tests are
continuous, nonnegative and nonzero. Division proves
the bound.
\end{proof}

The odd-cycle cancellation pertains to the Euclidean limit.
For example, group the six plaquettes at each raw base site
on the torus with \(a=1,\ell=2\). For one colour at unit
coupling, the third cumulant of their centred half-squares
at the three base sites \(0,e_4,e_3\) equals \(25/36864\).
To evaluate it directly, let \(A(z)\) be the matrix in
\eqref{gi:eq-lattice-curvature-fourier} before taking its
\(PQ\) entry. If \(b_S(v)=\mathbf 1_S\wedge v\), with
\(\mathbf 1_S\) the coordinate indicator vector of
\(S\subset\{1,2,3,4\}\), the two-site frequencies give
\[
 192A(z)=\sum_{S\ne\varnothing}
            \frac{12}{|S|}(-1)^{\sum_{\mu\in S}z_\mu} b_S b_S^{\mathsf T}.
\]
In the ordered plane basis \(12,13,14,23,24,34\) the diagonal
of \((192A(e_4))(192A(e_3+e_4))(192A(e_3))\) is
\((20096,6240,6240,6240,6240,-40256)\).
All six ordered cycles in the third cumulant have the same
trace, by cyclicity and transpose. Dividing the sum of these
diagonal entries by \(192^3\) gives the stated cumulant.

For the two-point estimate fix \(d>0\) and \(0<\vartheta\le1\).
Take nonzero, nonnegative Lipschitz tests \(f,h\), each
supported in a ball of radius \(d\), with
\begin{equation}\label{gi:eq-free-shape-class}
 \begin{gathered}
 d\le\operatorname{dist}(\operatorname{supp} f,
                            \operatorname{supp} h)\le6d,\\
 \int f\ge\vartheta\|f\|_\infty d^4,\qquad
 \operatorname{Lip}(f)\le\|f\|_\infty/(\vartheta d),
 \end{gathered}
\end{equation}
and the same last two inequalities for \(h\).
Amplitudes are arbitrary. Dividing by the amplitude and
dilating a support ball to radius one gives a profile with
integral at least \(\vartheta\) and Lipschitz constant at
most \(\vartheta^{-1}\). For \(d\le\ell/100\) these conditions
also define a class on the torus. To evaluate \(\mathcal K\)
there, lift the two support balls so a minimizing pair of
support points has the short displacement. Each ball lifts
injectively, and all resulting cross-displacements have
length at most \(10d\) by the two support diameters.
Their Euclidean lengths are the torus distances. The value
of \(\mathcal K\) is thus independent of the simultaneous
periodic translate of the chosen lifts.

To describe coarsening let \(L\) be the fixed odd block factor
and \(b=L^j a\) for a nonnegative integer \(j\).
For \(j\ge1\) use the compatible periodic hierarchy on which
each period contains an integer number of blocks at scale
\(b\) and at least two coarse sites.
For the corresponding projections \(P_j\),
Proposition~\ref{gi:prop-free-slices} identifies the
conditional mean curvature as the Gaussian with covariance
\(g^2P_j\).
Let \(Q_a^{[j]}\) be \eqref{gi:eq-free-lattice-density} for
that curvature, centred with respect to its own law.

\gineedspace{15\baselineskip}
\begin{theorem}[The uniform separated two-point comparison]
\label{gi:thm-free-separated-comparison}
For the shape class \eqref{gi:eq-free-shape-class} with
\(10a\le d\le\ell/100\), the full curvature covariance obeys
\[
 \left|
 \frac{\operatorname{Cov}(Q_a(f),Q_a(h))}
             {C_Gg^4\mathcal K(f,h)}-1
 \right|
 \le C\bigl[(a/d)^{1/2}+(d/\ell)^4+a/(\vartheta^3 d)\bigr].
\]
The denominator is strictly positive.
If \(b=L^j a\le d\) in the compatible hierarchy, then the
corresponding bound for \(Q_a^{[j]}\) is
\begin{equation}\label{gi:eq-free-background-error}
 \left|
 \frac{\operatorname{Cov}(Q_a^{[j]}(f),Q_a^{[j]}(h))}
             {C_Gg^4\mathcal K(f,h)}-1
 \right|
 \le C_L\bigl[(a/d)^{1/2}+(d/\ell)^4+a/(\vartheta^3 d)
                           +\mathbf 1_{\{j\ge1\}}e^{-c_Ld/b}\bigr].
\end{equation}
The constants are independent of depth, the torus side, mesh,
test amplitudes and positions, and \(g\).
\end{theorem}

\begin{proof}
The supports and mass lower bounds give
\[
        c\vartheta^2\|f\|_\infty\|h\|_\infty
        \le \mathcal K(f,h)\le C\|f\|_\infty\|h\|_\infty.
\]
For the upper bound use the volume of the two balls and
the lower separation. For the lower bound use the maximum
distance \(10d\) and the two mass bounds.

For \(j\ge1\), \cite[Proposition~6.4]{MB} gives, for
\(1\le i\le j\), the counting-coordinate entry bound
\[
  |(P_{i-1}-P_i)(p,q)|
     \le C_L(a/b_i)^4 e^{-c_L r/b_i},
                         \qquad b_i=L^{i-1}a .
\]
Here \(r\) may be taken as the barycentre distance, because
moving the chosen plaquette reference sites by \(O(a)\)
changes the exponential by a fixed factor.
Multiplying by \(a^{-4}\) and summing gives a physical kernel
error bounded by \(C_L r^{-4}e^{-c'_L r/b}\) for \(r\ge b\).
In fact \(r/b_i=L^{j-i+1}r/b\), and the sum of
\((r/b_i)^4 e^{-c_L r/b_i}\) is bounded by a smaller
exponential in \(r/b\). For \(j=0\) the difference vanishes.
This derives the coarsened kernel error by subtracting
the fluctuation projections from the full projection.

On the cross-support lattice pairs write \(E\) for the
sum of the mesh and volume errors and, for the coarsened
law, the exponential error on the right of
\eqref{gi:eq-free-background-error}, without the quadrature
term. Each physical entry of the relevant covariance
divided by \(g^2\) differs from
\(\mathsf K_{PQ}(u)\) by at most \(C_L E|u|^{-4}\).
Its squared entry therefore differs by at most
\(C_L E|u|^{-8}\), since \(E\) is bounded in the stated
range. Positivity of the tests bounds the summed error by
this factor times the corresponding Riemann sum of
\(f(x)h(y)|u|^{-8}\).

For each fixed pair of plane labels, the shifted-grid
quadrature for \(f(x)h(y)|u|^{-8}\) has absolute error at
most \(C a\|f\|_\infty\|h\|_\infty/(\vartheta d)\).
Indeed the mesh cells meeting the product of the supports
are contained in the product of two balls of radius
\(d+2a\), and throughout them the cross-distance is
bounded below by \(d-4a\). The product rule bounds the
Lipschitz norm of the integrand there by
\(C\vartheta^{-1}d^{-9}\|f\|_\infty\|h\|_\infty\).
Using the cell diameter \(Ca\) and the product volume
\(Cd^8\) proves the bound. The same argument applies to
the integrand with \(\mathsf K_{PQ}(u)^2\) in place of
\(|u|^{-8}\), by Lemma~\ref{gi:lem-curvature-reflection}
and its differentiated expression.
The relative quadrature cost is at most
\(Ca/(\vartheta^3d)\) by the lower bound on \(\mathcal K\).
Multiplication of that cost by the bounded \(E\) still
gives a term of the stated order.

Finally sum over the planes and apply
Lemma~\ref{gi:lem-curvature-reflection} to the continuum
integral. The Wick factor \(d_Gg^4/2\) in
\eqref{gi:eq-free-lattice-wick} then gives the coefficient
in \eqref{gi:eq-free-coefficient}.
\end{proof}

\clearpage
\begin{thebibliography}{9}
\raggedright
\bibitem{BA} \emph{Block averaging and gauge geometry for four-dimensional lattice Yang--Mills theory}, companion manuscript, 2026.
\bibitem{MB} \emph{Linear lattice PDE with block constraints}, companion manuscript, 2026.
\bibitem{Bal84a} T. Ba\l aban, \emph{Propagators and renormalization transformations for lattice gauge theories. I}, Comm. Math. Phys. \textbf{95} (1984), 17--40. \href{https://doi.org/10.1007/BF01215753}{doi:10.1007/BF01215753}.
\bibitem{Bal84b} T. Ba\l aban, \emph{Propagators and renormalization transformations for lattice gauge theories. II}, Comm. Math. Phys. \textbf{96} (1984), 223--250. \href{https://doi.org/10.1007/BF01240221}{doi:10.1007/BF01240221}.
\bibitem{Bal85p} T. Ba\l aban, \emph{Propagators for lattice gauge theories in a background field}, Comm. Math. Phys. \textbf{99} (1985), 389--434. \href{https://doi.org/10.1007/BF01240355}{doi:10.1007/BF01240355}.
\end{thebibliography}
\end{document}
